% Classroom adaptation of the six core General Worked Problems, in source order.
% Advanced Problems are intentionally excluded. Each family has 1 worked,
% 1 guided and 2 independent problems. All new distributions were solved
% by exact rational enumeration; full matching solutions are included below.
\makeatletter
\def\input@path{{../../../preamble/}}
\makeatother
\input{preamble.tex}
\input{graphs/highlights.tex}
% \GTenHighlightsfalse % Uncomment to disable instructional highlights.
\tikzset{gten marker/.append style={inner xsep=.6mm,inner ysep=.3mm}}
% Local C2 roles follow C1: quantities orange; probabilities in related blues.
\colorlet{cTwoMassMarker}{GTenHLProbability}
\definecolor{cTwoCDFMarker}{HTML}{43A5C9}
\newcommand{\cTwoValue}[1]{\HighlightStroke{GTenHLValue}{#1}}
\newcommand{\cTwoOutcome}[1]{\HighlightStroke{GTenHLOutcome}{#1}}
\newcommand{\cTwoMassHL}[1]{\HighlightStroke{cTwoMassMarker}{#1}}
\newcommand{\cTwoCDFHL}[1]{\HighlightStroke{cTwoCDFMarker}{#1}}
\colorlet{cTwoMassColor}{blue!75!black}
\colorlet{cTwoCDFColor}{cyan!65!blue!80!black}
\colorlet{cTwoGuideColor}{cTwoMassColor!35}
\ifGTenHighlights\else
  \colorlet{cTwoMassColor}{black}
  \colorlet{cTwoCDFColor}{black}
  \colorlet{cTwoGuideColor}{gray}
\fi
\input{graphs/G10_T2_bar_graphs.tex}
\usepackage{needspace}
\usepackage{varwidth}
\geometry{letterpaper,margin=0.25in}
\pagestyle{empty}
\linespread{1.04}
\setlength{\parskip}{4pt}
\setlength{\tabcolsep}{6pt}
\setlength{\extrarowheight}{2pt}
\setlength{\emergencystretch}{2em}
\hypersetup{pdftitle={Grade 10 C2 Classroom Activity: PMF and CDF Tables and Graphs}}
\newif\ifprintsolutions
\printsolutionsfalse % Set \printsolutionstrue to append the complete review solutions.
\newcolumntype{Y}{>{\centering\arraybackslash}X}
\newcommand{\stageheading}[1]{\par\Needspace{6\baselineskip}\addvspace{.65\baselineskip}{\centering\large\bfseries #1\par}\nopagebreak[4]}
\newcommand{\problemtitle}[3]{\par\Needspace{4\baselineskip}\addvspace{.4\baselineskip}\noindent{\setlength{\fboxsep}{2pt}\colorbox{black!10}{\begin{varwidth}[t]{.75\linewidth}\bfseries\strut #1. #2\end{varwidth}}}\hfill\textbf{Type #3}\par\nopagebreak[4]}
\newcommand{\task}[2]{\par\noindent\textbf{Task #1 --- \ifnum#1=1 PMF table and graph\else CDF table and graph\fi.} #2\par}
\newcommand{\responseline}{\par\noindent\rule{0pt}{1.55em}\hrulefill\par}
\newcommand{\worklines}[1]{\begingroup\setlength{\parskip}{0pt}\foreach \i in {1,...,#1}{\responseline}\endgroup}
\newcommand{\blankgraph}[1]{\begin{tikzpicture}\begin{axis}[width=\linewidth,height=4.4cm,axis lines=left,xmin=0,xmax=10,ymin=0,ymax=1.05,xtick=\empty,ytick={0,.25,.5,.75,1},grid=major,grid style={gray!20},xlabel={$x$ (USD)},ylabel={#1},tick label style={font=\small},label style={font=\small}]\end{axis}\end{tikzpicture}}
\newcommand{\blanktable}[1]{\begin{tabularx}{\linewidth}{|Y|Y|}\hline #1\\\hline\rule{0pt}{.6cm}&\\\hline\rule{0pt}{.6cm}&\\\hline\rule{0pt}{.6cm}&\\\hline\rule{0pt}{.6cm}&\\\hline\rule{0pt}{.6cm}&\\\hline\rule{0pt}{.6cm}&\\\hline\end{tabularx}}
\pgfplotsset{classCDF/.style={width=\dimexpr\linewidth-6pt\relax,height=4.15cm,axis lines=left,ymin=-.06,ymax=1.1,ytick={0,.25,.5,.75,1},ylabel={$\cTwoCDFHL{F_X(x)}$},xlabel={$\cTwoValue{x}$ (\cTwoValue{USD})},grid=major,grid style={gray!15},tick label style={font=\small},label style={font=\small},clip=false}}
\begin{document}
{\centering{\Large\bfseries C2 -- Constructing PMF and CDF Tables and Graphs\par}
Global Economics \textbar{} Grade 10 \textbar{} Term 2 \textbar{} Lesson 1\par}
\textbf{C2:} Interprets PMFs and CDFs by verifying their defining properties and using them to calculate point, interval, cumulative, and complementary probabilities in financial contexts.

\stageheading{Equations and basic procedure}
\noindent\begin{tabularx}{\linewidth}{@{}l l X@{}}\toprule
Symbol or notation & Concept & Brief meaning\\\midrule
$\cTwoValue{X},\cTwoValue{x},\cTwoValue{S_X}$ & \cTwoValue{Quantity} & Variable, value or cutoff, and support; orange.\\
$\cTwoOutcome{\omega}$ & \cTwoOutcome{Complete outcome} & One complete result; violet.\\
$\cTwoMassHL{p_X(x)}$ & \cTwoMassHL{PMF} & Probability at one value; blue bars or stems.\\
$\cTwoCDFHL{F_X(x)}$ & \cTwoCDFHL{CDF} & Probability through a cutoff; related blue steps.\\\bottomrule
\end{tabularx}

All data are hypothetical. Monetary values are in USD; probabilities have no units. This activity constructs and interprets the four representations used in C2. Keep exact fractions or terminating decimals.

\task{1}{Define the variable and its units. Translate complete outcomes into values. Multiply probabilities within independent ordered outcomes; add mutually exclusive outcomes sharing a value. Recover any missing probability using total probability 1. Apply caps, discounts and costs before grouping. List the sorted support, make the PMF table, and draw separated bars or stems at the actual numerical values. Check every mass and the total.}
\task{2}{Accumulate the PMF from left to right. Give a CDF table covering all real intervals, including $x<\min S_X$ with height 0 and $x\geq\max S_X$ with height 1. Draw horizontal steps: at each jump, the upper filled point is included and the lower open point is excluded. Check nondecreasing heights, right-continuity, tail limits 0 and 1, and jumps equal to PMF masses. Interpret a cumulative height in context.}
\textbf{Graph checks.} Use numerical spacing, including unequal gaps and negative results. PMF bars represent individual values, not intervals; bar width and area play no role. Do not connect PMF values or CDF jumps with sloping lines. Draw arrows for both CDF tails. A negative monetary value can have a positive probability.

\textbf{Family route (kept in this order in every stage).}
\begin{center}\begin{tabularx}{.96\linewidth}{lX}\toprule
Type & Structure to preserve \\\midrule
A: Payment fees & Four methods map to three fees; group both card methods.\\
B: Late fees & Four delay categories map to three fees; group both shorter delays.\\
C: Uneven buyback payments & Six material categories map to four uneven values; group two pairs.\\
D: Missing probability and capped payout & Recover a claim-count probability, apply the payout cap, then group.\\
E: Independent purchases and discount & List eight ordered decisions, multiply, apply the all-buy discount, then group.\\
F: Signed two-stage net results & List six independent ordered pairs, subtract one cost, group break-even results, sort.\\\bottomrule
\end{tabularx}\end{center}
\textbf{Student instructions.} Complete both tasks for every problem. Blank table rows may be left unused. Label a suitable horizontal scale on each graph; its printed grid does not prescribe the support. Use the calculation area for outcome mappings, checks and interpretations.
\vfill
Name: \hrulefill\quad Class: \rule{2cm}{.4pt}\quad Date: \rule{2.5cm}{.4pt}
\clearpage
\stageheading{Worked example}
\problemtitle{1}{Museum kiosk payments}{A}
\textbf{Payment fees.} A museum kiosk receives one payment. The only methods are cash, debit card, credit card and mobile wallet, with probabilities .15, .25, .35, .25, respectively. Cash has no fee; either card costs 2 USD; a mobile wallet costs 4 USD. There are no other charges. Record the processing fee.
\task{1}{Let $\cTwoValue{X}$ be the \cTwoValue{processing fee} for one payment, in \cTwoValue{USD}. Cash maps to \cTwoValue{0}, either card to \cTwoValue{2}, and wallet to \cTwoValue{4}. The two card events are mutually exclusive, so their masses add: $\cTwoMassHL{p_X(2)}=\cTwoMassHL{.25}+\cTwoMassHL{.35}=\cTwoMassHL{0.6}$.}
\begin{center}
\begin{tabular}{ccc}\toprule
Complete outcome & Probability & $\cTwoValue{X}$ (\cTwoValue{USD}) \\ \midrule
\cTwoOutcome{Cash} & $\cTwoMassHL{0.15}$ & $\cTwoValue{0}$ \\
\cTwoOutcome{Debit card} & $\cTwoMassHL{0.25}$ & $\cTwoValue{2}$ \\
\cTwoOutcome{Credit card} & $\cTwoMassHL{0.35}$ & $\cTwoValue{2}$ \\
\cTwoOutcome{Wallet} & $\cTwoMassHL{0.25}$ & $\cTwoValue{4}$ \\
\bottomrule\end{tabular}
\end{center}
$\cTwoValue{S_X}=\cTwoValue{\{0,2,4\}}$; $p_X(x)=0$ elsewhere. The grouped PMF is:
\par\noindent\begin{minipage}[t]{.37\linewidth}\vspace{0pt}\centering
\begin{tabular}{cc}\toprule
$\cTwoValue{x}$ (\cTwoValue{USD}) & $\cTwoMassHL{p_X(x)}$ \\ \midrule
$\cTwoValue{0}$ & $\cTwoMassHL{0.15}$ \\
$\cTwoValue{2}$ & $\cTwoMassHL{0.6}$ \\
$\cTwoValue{4}$ & $\cTwoMassHL{0.25}$ \\
\bottomrule\end{tabular}
\end{minipage}\hfill\begin{minipage}[t]{.60\linewidth}\vspace{0pt}\centering
\GTenPMFGraph[color=cTwoMassColor!20,xhighlight=GTenHLValue,yhighlight=cTwoMassMarker,width=\linewidth,height=4.15cm,xlabel={$\cTwoValue{x}$ (\cTwoValue{USD})},ylabel={$\cTwoMassHL{p_X(x)}$},xstep=1]{0,2,4}{0.15,0.6,0.25}
\end{minipage}\par
\textbf{PMF check.} All masses are positive and $0.15+0.6+0.25=1$. The graph uses the actual numerical spacing; gaps carry no mass.
\task{2}{Accumulate in increasing support order: $\cTwoCDFHL{0.15};\ \cTwoCDFHL{0.75};\ \cTwoCDFHL{1}$. Each new height is the previous height plus the next PMF mass.}
\par\noindent\begin{minipage}[t]{.37\linewidth}\vspace{0pt}\centering
\begin{tabular}{cc}\toprule
Cutoff interval & $\cTwoCDFHL{F_X(x)}$ \\ \midrule
$x<0$ & $\cTwoCDFHL{0}$ \\
$0\leq x<2$ & $\cTwoCDFHL{0.15}$ \\
$2\leq x<4$ & $\cTwoCDFHL{0.75}$ \\
$x\geq 4$ & $\cTwoCDFHL{1}$ \\
\bottomrule\end{tabular}
\end{minipage}\hfill\begin{minipage}[t]{.60\linewidth}\vspace{0pt}\centering
\begin{tikzpicture}\begin{axis}[classCDF,xmin=-2,xmax=6,xtick={0,2,4}]
\addplot[cTwoCDFColor,thick,<-] coordinates {(-2,0)(0,0)};
\addplot[cTwoCDFColor,thick] coordinates {(0,0.15)(2,0.15)};
\addplot[cTwoCDFColor,thick] coordinates {(2,0.75)(4,0.75)};
\addplot[cTwoCDFColor,thick,->] coordinates {(4,1)(6,1)};
\addplot[only marks,mark=*,cTwoCDFColor,mark size=1.8pt,mark options={fill=white}] coordinates {(0,0)(2,0.15)(4,0.75)};
\addplot[only marks,mark=*,cTwoCDFColor,mark size=1.8pt] coordinates {(0,0.15)(2,0.75)(4,1)};
\end{axis}\end{tikzpicture}
\end{minipage}\par
\textbf{CDF check.} All heights lie in $[0,1]$ and never decrease. Filled left/open right endpoints make the graph right-continuous. Horizontal tails give limits 0 and 1; the CDF reaches 1 at 4 USD. The jumps, $0.15,\ 0.6,\ 0.25$, reproduce the PMF.
\textbf{Interpretation.} $F_X(2)=0.75$: the probability that the fee is at most 2 USD (cash or either card) is 0.75.
\vfill
\clearpage
\problemtitle{2}{Library equipment returns}{B}
\textbf{Late fees.} A rental service records the return of one item of library equipment. The exhaustive categories are on time or early, 1--5 whole minutes late, 6--15 whole minutes late, and more than 15 minutes late, with probabilities .20, .30, .15, .35, respectively. The late fee is 0 USD, 2 USD, 2 USD and 4 USD, respectively. No other charges apply. Record the late fee.
\task{1}{Let $\cTwoValue{X}$ be the late fee for one return, in USD. The two shorter delays both map to 2 USD. Their mutually exclusive masses add: $p_X(2)=.30+.15=0.45$. On-time returns map to 0 and the longest delay maps to 4.}
\begin{center}
\begin{tabular}{ccc}\toprule
Complete outcome & Probability & $\cTwoValue{X}$ (\cTwoValue{USD}) \\ \midrule
\cTwoOutcome{On time/early} & $\cTwoMassHL{0.2}$ & $\cTwoValue{0}$ \\
\cTwoOutcome{1--5 min} & $\cTwoMassHL{0.3}$ & $\cTwoValue{2}$ \\
\cTwoOutcome{6--15 min} & $\cTwoMassHL{0.15}$ & $\cTwoValue{2}$ \\
\cTwoOutcome{Over 15 min} & $\cTwoMassHL{0.35}$ & $\cTwoValue{4}$ \\
\bottomrule\end{tabular}
\end{center}
$\cTwoValue{S_X}=\cTwoValue{\{0,2,4\}}$; $p_X(x)=0$ elsewhere. The grouped PMF is:
\par\noindent\begin{minipage}[t]{.37\linewidth}\vspace{0pt}\centering
\begin{tabular}{cc}\toprule
$\cTwoValue{x}$ (\cTwoValue{USD}) & $\cTwoMassHL{p_X(x)}$ \\ \midrule
$\cTwoValue{0}$ & $\cTwoMassHL{0.2}$ \\
$\cTwoValue{2}$ & $\cTwoMassHL{0.45}$ \\
$\cTwoValue{4}$ & $\cTwoMassHL{0.35}$ \\
\bottomrule\end{tabular}
\end{minipage}\hfill\begin{minipage}[t]{.60\linewidth}\vspace{0pt}\centering
\GTenPMFGraph[color=cTwoMassColor!20,xhighlight=GTenHLValue,yhighlight=cTwoMassMarker,width=\linewidth,height=4.15cm,xlabel={$\cTwoValue{x}$ (\cTwoValue{USD})},ylabel={$\cTwoMassHL{p_X(x)}$},xstep=1]{0,2,4}{0.2,0.45,0.35}
\end{minipage}\par
\textbf{PMF check.} All masses are positive and $0.2+0.45+0.35=1$. The graph uses the actual numerical spacing; gaps carry no mass.
\task{2}{Accumulate in increasing support order: $0.2;\ 0.65;\ 1$. Each new height is the previous height plus the next PMF mass.}
\par\noindent\begin{minipage}[t]{.37\linewidth}\vspace{0pt}\centering
\begin{tabular}{cc}\toprule
Cutoff interval & $\cTwoCDFHL{F_X(x)}$ \\ \midrule
$x<0$ & $\cTwoCDFHL{0}$ \\
$0\leq x<2$ & $\cTwoCDFHL{0.2}$ \\
$2\leq x<4$ & $\cTwoCDFHL{0.65}$ \\
$x\geq 4$ & $\cTwoCDFHL{1}$ \\
\bottomrule\end{tabular}
\end{minipage}\hfill\begin{minipage}[t]{.60\linewidth}\vspace{0pt}\centering
\begin{tikzpicture}\begin{axis}[classCDF,xmin=-2,xmax=6,xtick={0,2,4}]
\addplot[cTwoCDFColor,thick,<-] coordinates {(-2,0)(0,0)};
\addplot[cTwoCDFColor,thick] coordinates {(0,0.2)(2,0.2)};
\addplot[cTwoCDFColor,thick] coordinates {(2,0.65)(4,0.65)};
\addplot[cTwoCDFColor,thick,->] coordinates {(4,1)(6,1)};
\addplot[only marks,mark=*,cTwoCDFColor,mark size=1.8pt,mark options={fill=white}] coordinates {(0,0)(2,0.2)(4,0.65)};
\addplot[only marks,mark=*,cTwoCDFColor,mark size=1.8pt] coordinates {(0,0.2)(2,0.65)(4,1)};
\end{axis}\end{tikzpicture}
\end{minipage}\par
\textbf{CDF check.} All heights lie in $[0,1]$ and never decrease. Filled left/open right endpoints make the graph right-continuous. Horizontal tails give limits 0 and 1; the CDF reaches 1 at 4 USD. The jumps, $0.2,\ 0.45,\ 0.35$, reproduce the PMF.
\textbf{Interpretation.} $F_X(2)=0.65$: the probability that the late fee is at most 2 USD (the return is no more than 15 minutes late) is 0.65.
\vfill
\clearpage
\problemtitle{3}{Craft-material buyback}{C}
\textbf{Uneven buyback payments.} A craft-material centre buys one standard-sized bag containing only one material type. The exhaustive types are rejected material, paper, cardboard, plastic, steel and aluminium, with probabilities .15, .10, .20, .25, .10, .20, respectively. Payments are 0 USD for rejected material, 3 USD for either paper or cardboard, 7 USD for either plastic or steel, and 12 USD for aluminium. There are no other payments. Record the buyback payment.
\task{1}{Let $\cTwoValue{X}$ be the buyback payment for one bag, in USD. Paper and cardboard share 3 USD, so their probabilities add to 0.3. Plastic and steel share 7 USD, giving 0.35. Keep the actual gaps 3, 4 and 5 USD on both graphs.}
\begin{center}
\begin{tabular}{ccc}\toprule
Complete outcome & Probability & $\cTwoValue{X}$ (\cTwoValue{USD}) \\ \midrule
\cTwoOutcome{Rejected} & $\cTwoMassHL{0.15}$ & $\cTwoValue{0}$ \\
\cTwoOutcome{Paper} & $\cTwoMassHL{0.1}$ & $\cTwoValue{3}$ \\
\cTwoOutcome{Cardboard} & $\cTwoMassHL{0.2}$ & $\cTwoValue{3}$ \\
\cTwoOutcome{Plastic} & $\cTwoMassHL{0.25}$ & $\cTwoValue{7}$ \\
\cTwoOutcome{Steel} & $\cTwoMassHL{0.1}$ & $\cTwoValue{7}$ \\
\cTwoOutcome{Aluminium} & $\cTwoMassHL{0.2}$ & $\cTwoValue{12}$ \\
\bottomrule\end{tabular}
\end{center}
$\cTwoValue{S_X}=\cTwoValue{\{0,3,7,12\}}$; $p_X(x)=0$ elsewhere. The grouped PMF is:
\par\noindent\begin{minipage}[t]{.37\linewidth}\vspace{0pt}\centering
\begin{tabular}{cc}\toprule
$\cTwoValue{x}$ (\cTwoValue{USD}) & $\cTwoMassHL{p_X(x)}$ \\ \midrule
$\cTwoValue{0}$ & $\cTwoMassHL{0.15}$ \\
$\cTwoValue{3}$ & $\cTwoMassHL{0.3}$ \\
$\cTwoValue{7}$ & $\cTwoMassHL{0.35}$ \\
$\cTwoValue{12}$ & $\cTwoMassHL{0.2}$ \\
\bottomrule\end{tabular}
\end{minipage}\hfill\begin{minipage}[t]{.60\linewidth}\vspace{0pt}\centering
\GTenPMFGraph[color=cTwoMassColor!20,xhighlight=GTenHLValue,yhighlight=cTwoMassMarker,width=\linewidth,height=4.15cm,xlabel={$\cTwoValue{x}$ (\cTwoValue{USD})},ylabel={$\cTwoMassHL{p_X(x)}$},xstep=1]{0,3,7,12}{0.15,0.3,0.35,0.2}
\end{minipage}\par
\textbf{PMF check.} All masses are positive and $0.15+0.3+0.35+0.2=1$. The graph uses the actual numerical spacing; gaps carry no mass.
\task{2}{Accumulate in increasing support order: $0.15;\ 0.45;\ 0.8;\ 1$. Each new height is the previous height plus the next PMF mass.}
\par\noindent\begin{minipage}[t]{.37\linewidth}\vspace{0pt}\centering
\begin{tabular}{cc}\toprule
Cutoff interval & $\cTwoCDFHL{F_X(x)}$ \\ \midrule
$x<0$ & $\cTwoCDFHL{0}$ \\
$0\leq x<3$ & $\cTwoCDFHL{0.15}$ \\
$3\leq x<7$ & $\cTwoCDFHL{0.45}$ \\
$7\leq x<12$ & $\cTwoCDFHL{0.8}$ \\
$x\geq 12$ & $\cTwoCDFHL{1}$ \\
\bottomrule\end{tabular}
\end{minipage}\hfill\begin{minipage}[t]{.60\linewidth}\vspace{0pt}\centering
\begin{tikzpicture}\begin{axis}[classCDF,xmin=-3,xmax=15,xtick={0,3,7,12}]
\addplot[cTwoCDFColor,thick,<-] coordinates {(-3,0)(0,0)};
\addplot[cTwoCDFColor,thick] coordinates {(0,0.15)(3,0.15)};
\addplot[cTwoCDFColor,thick] coordinates {(3,0.45)(7,0.45)};
\addplot[cTwoCDFColor,thick] coordinates {(7,0.8)(12,0.8)};
\addplot[cTwoCDFColor,thick,->] coordinates {(12,1)(15,1)};
\addplot[only marks,mark=*,cTwoCDFColor,mark size=1.8pt,mark options={fill=white}] coordinates {(0,0)(3,0.15)(7,0.45)(12,0.8)};
\addplot[only marks,mark=*,cTwoCDFColor,mark size=1.8pt] coordinates {(0,0.15)(3,0.45)(7,0.8)(12,1)};
\end{axis}\end{tikzpicture}
\end{minipage}\par
\textbf{CDF check.} All heights lie in $[0,1]$ and never decrease. Filled left/open right endpoints make the graph right-continuous. Horizontal tails give limits 0 and 1; the CDF reaches 1 at 12 USD. The jumps, $0.15,\ 0.3,\ 0.35,\ 0.2$, reproduce the PMF.
\textbf{Interpretation.} $F_X(7)=0.8$: the probability that the buyback payment is at most 7 USD (every type except aluminium) is 0.8.
\vfill
\clearpage
\problemtitle{4}{Parcel damage reimbursements}{D}
\textbf{Missing probability and capped payout.} A parcel insurer can receive only 0, 1, 2, 3 or 4 claims in one day. Their probabilities, in that order, are .20, .25, a missing probability, .18, .12. It pays 3 USD per claim, capped at 9 USD in total for the day; it pays the smaller of 3 times the count and 9 USD. There are no other payments. Record the daily payout.
\task{1}{Let $\cTwoValue{X}$ be the total payout for one day, in USD. The missing probability is $1-(.20+.25+.18+.12)=.25$. Apply $X=\min(3N,9)$ to the claim count $N$: payouts are $0,3,6,9,9$. Four claims would pay 12 USD without the cap. Group the last two categories.}
\begin{center}
\begin{tabular}{ccc}\toprule
Complete outcome & Probability & $\cTwoValue{X}$ (\cTwoValue{USD}) \\ \midrule
\cTwoOutcome{0 claims} & $\cTwoMassHL{0.2}$ & $\cTwoValue{0}$ \\
\cTwoOutcome{1 claims} & $\cTwoMassHL{0.25}$ & $\cTwoValue{3}$ \\
\cTwoOutcome{2 claims} & $\cTwoMassHL{0.25}$ & $\cTwoValue{6}$ \\
\cTwoOutcome{3 claims} & $\cTwoMassHL{0.18}$ & $\cTwoValue{9}$ \\
\cTwoOutcome{4 claims} & $\cTwoMassHL{0.12}$ & $\cTwoValue{9}$ \\
\bottomrule\end{tabular}
\end{center}
$\cTwoValue{S_X}=\cTwoValue{\{0,3,6,9\}}$; $p_X(x)=0$ elsewhere. The grouped PMF is:
\par\noindent\begin{minipage}[t]{.37\linewidth}\vspace{0pt}\centering
\begin{tabular}{cc}\toprule
$\cTwoValue{x}$ (\cTwoValue{USD}) & $\cTwoMassHL{p_X(x)}$ \\ \midrule
$\cTwoValue{0}$ & $\cTwoMassHL{0.2}$ \\
$\cTwoValue{3}$ & $\cTwoMassHL{0.25}$ \\
$\cTwoValue{6}$ & $\cTwoMassHL{0.25}$ \\
$\cTwoValue{9}$ & $\cTwoMassHL{0.3}$ \\
\bottomrule\end{tabular}
\end{minipage}\hfill\begin{minipage}[t]{.60\linewidth}\vspace{0pt}\centering
\GTenPMFGraph[color=cTwoMassColor!20,xhighlight=GTenHLValue,yhighlight=cTwoMassMarker,width=\linewidth,height=4.15cm,xlabel={$\cTwoValue{x}$ (\cTwoValue{USD})},ylabel={$\cTwoMassHL{p_X(x)}$},xstep=1]{0,3,6,9}{0.2,0.25,0.25,0.3}
\end{minipage}\par
\textbf{PMF check.} All masses are positive and $0.2+0.25+0.25+0.3=1$. The graph uses the actual numerical spacing; gaps carry no mass.
\task{2}{Accumulate in increasing support order: $0.2;\ 0.45;\ 0.7;\ 1$. Each new height is the previous height plus the next PMF mass.}
\par\noindent\begin{minipage}[t]{.37\linewidth}\vspace{0pt}\centering
\begin{tabular}{cc}\toprule
Cutoff interval & $\cTwoCDFHL{F_X(x)}$ \\ \midrule
$x<0$ & $\cTwoCDFHL{0}$ \\
$0\leq x<3$ & $\cTwoCDFHL{0.2}$ \\
$3\leq x<6$ & $\cTwoCDFHL{0.45}$ \\
$6\leq x<9$ & $\cTwoCDFHL{0.7}$ \\
$x\geq 9$ & $\cTwoCDFHL{1}$ \\
\bottomrule\end{tabular}
\end{minipage}\hfill\begin{minipage}[t]{.60\linewidth}\vspace{0pt}\centering
\begin{tikzpicture}\begin{axis}[classCDF,xmin=-3,xmax=12,xtick={0,3,6,9}]
\addplot[cTwoCDFColor,thick,<-] coordinates {(-3,0)(0,0)};
\addplot[cTwoCDFColor,thick] coordinates {(0,0.2)(3,0.2)};
\addplot[cTwoCDFColor,thick] coordinates {(3,0.45)(6,0.45)};
\addplot[cTwoCDFColor,thick] coordinates {(6,0.7)(9,0.7)};
\addplot[cTwoCDFColor,thick,->] coordinates {(9,1)(12,1)};
\addplot[only marks,mark=*,cTwoCDFColor,mark size=1.8pt,mark options={fill=white}] coordinates {(0,0)(3,0.2)(6,0.45)(9,0.7)};
\addplot[only marks,mark=*,cTwoCDFColor,mark size=1.8pt] coordinates {(0,0.2)(3,0.45)(6,0.7)(9,1)};
\end{axis}\end{tikzpicture}
\end{minipage}\par
\textbf{CDF check.} All heights lie in $[0,1]$ and never decrease. Filled left/open right endpoints make the graph right-continuous. Horizontal tails give limits 0 and 1; the CDF reaches 1 at 9 USD. The jumps, $0.2,\ 0.25,\ 0.25,\ 0.3$, reproduce the PMF.
\textbf{Interpretation.} $F_X(6)=0.7$: the probability that the payout is at most 6 USD (at most two claims) is 0.7.
\vfill
\clearpage
\problemtitle{5}{Badge sales}{E}
\textbf{Independent purchases and discount.} A school enterprise offers one badge to each of three visitors. Independently, each buys one with probability 0.4 or passes with probability 0.6, regardless of the discount. Each item costs 4 USD. Only if all three buy, subtract a 2 USD discount from the combined bill. Otherwise there is no discount. There are no other sales or charges. Record total revenue.
\task{1}{Let $\cTwoValue{X}$ be the total revenue from the three visitors, in USD. Use B for buying and N for passing, in visitor order. Multiply within each complete sequence: e.g. $P(\mathrm{BNN})=0.4(0.6)^2=0.144$. The three one-buyer and three two-buyer outcomes each form a group. Only BBB receives the discount: $3(4)-2=10$ USD.}
\begin{center}
\begin{tabular}{ccc}\toprule
Complete outcome & Probability & $\cTwoValue{X}$ (\cTwoValue{USD}) \\ \midrule
\cTwoOutcome{NNN} & $\cTwoMassHL{0.216}$ & $\cTwoValue{0}$ \\
\cTwoOutcome{BNN} & $\cTwoMassHL{0.144}$ & $\cTwoValue{4}$ \\
\cTwoOutcome{NBN} & $\cTwoMassHL{0.144}$ & $\cTwoValue{4}$ \\
\cTwoOutcome{NNB} & $\cTwoMassHL{0.144}$ & $\cTwoValue{4}$ \\
\cTwoOutcome{BBN} & $\cTwoMassHL{0.096}$ & $\cTwoValue{8}$ \\
\cTwoOutcome{BNB} & $\cTwoMassHL{0.096}$ & $\cTwoValue{8}$ \\
\cTwoOutcome{NBB} & $\cTwoMassHL{0.096}$ & $\cTwoValue{8}$ \\
\cTwoOutcome{BBB} & $\cTwoMassHL{0.064}$ & $\cTwoValue{10}$ \\
\bottomrule\end{tabular}
\end{center}
$\cTwoValue{S_X}=\cTwoValue{\{0,4,8,10\}}$; $p_X(x)=0$ elsewhere. The grouped PMF is:
\par\noindent\begin{minipage}[t]{.37\linewidth}\vspace{0pt}\centering
\begin{tabular}{cc}\toprule
$\cTwoValue{x}$ (\cTwoValue{USD}) & $\cTwoMassHL{p_X(x)}$ \\ \midrule
$\cTwoValue{0}$ & $\cTwoMassHL{0.216}$ \\
$\cTwoValue{4}$ & $\cTwoMassHL{0.432}$ \\
$\cTwoValue{8}$ & $\cTwoMassHL{0.288}$ \\
$\cTwoValue{10}$ & $\cTwoMassHL{0.064}$ \\
\bottomrule\end{tabular}
\end{minipage}\hfill\begin{minipage}[t]{.60\linewidth}\vspace{0pt}\centering
\GTenPMFGraph[color=cTwoMassColor!20,xhighlight=GTenHLValue,yhighlight=cTwoMassMarker,width=\linewidth,height=4.15cm,xlabel={$\cTwoValue{x}$ (\cTwoValue{USD})},ylabel={$\cTwoMassHL{p_X(x)}$},xstep=1]{0,4,8,10}{0.216,0.432,0.288,0.064}
\end{minipage}\par
\textbf{PMF check.} All masses are positive and $0.216+0.432+0.288+0.064=1$. The graph uses the actual numerical spacing; gaps carry no mass.
\task{2}{Accumulate in increasing support order: $0.216;\ 0.648;\ 0.936;\ 1$. Each new height is the previous height plus the next PMF mass.}
\par\noindent\begin{minipage}[t]{.37\linewidth}\vspace{0pt}\centering
\begin{tabular}{cc}\toprule
Cutoff interval & $\cTwoCDFHL{F_X(x)}$ \\ \midrule
$x<0$ & $\cTwoCDFHL{0}$ \\
$0\leq x<4$ & $\cTwoCDFHL{0.216}$ \\
$4\leq x<8$ & $\cTwoCDFHL{0.648}$ \\
$8\leq x<10$ & $\cTwoCDFHL{0.936}$ \\
$x\geq 10$ & $\cTwoCDFHL{1}$ \\
\bottomrule\end{tabular}
\end{minipage}\hfill\begin{minipage}[t]{.60\linewidth}\vspace{0pt}\centering
\begin{tikzpicture}\begin{axis}[classCDF,xmin=-2,xmax=12,xtick={0,4,8,10}]
\addplot[cTwoCDFColor,thick,<-] coordinates {(-2,0)(0,0)};
\addplot[cTwoCDFColor,thick] coordinates {(0,0.216)(4,0.216)};
\addplot[cTwoCDFColor,thick] coordinates {(4,0.648)(8,0.648)};
\addplot[cTwoCDFColor,thick] coordinates {(8,0.936)(10,0.936)};
\addplot[cTwoCDFColor,thick,->] coordinates {(10,1)(12,1)};
\addplot[only marks,mark=*,cTwoCDFColor,mark size=1.8pt,mark options={fill=white}] coordinates {(0,0)(4,0.216)(8,0.648)(10,0.936)};
\addplot[only marks,mark=*,cTwoCDFColor,mark size=1.8pt] coordinates {(0,0.216)(4,0.648)(8,0.936)(10,1)};
\end{axis}\end{tikzpicture}
\end{minipage}\par
\textbf{CDF check.} All heights lie in $[0,1]$ and never decrease. Filled left/open right endpoints make the graph right-continuous. Horizontal tails give limits 0 and 1; the CDF reaches 1 at 10 USD. The jumps, $0.216,\ 0.432,\ 0.288,\ 0.064$, reproduce the PMF.
\textbf{Interpretation.} $F_X(8)=0.936$: the probability that revenue is at most 8 USD (not all three visitors buy) is 0.936.
\vfill
\clearpage
\problemtitle{6}{Two-stage arts fund}{F}
\textbf{Signed two-stage net results.} A hypothetical classroom fund receives two cash returns. The first is 14 USD, 0 USD or 6 USD with probabilities .25, .35, .40, respectively. Independently, the second is 0 USD with probability 0.6 or 6 USD with probability 0.4. These are all possible returns. A fixed participation cost of 6 USD is charged once. There are no other receipts or charges. Record the sum of both returns minus the cost; a negative net result means a loss.
\task{1}{Let $\cTwoValue{X}$ be the signed net result after both stages and the cost, in USD. For each ordered pair, multiply independent probabilities, then compute $X=R_1+R_2-6$. For example, $(14,0)$ has probability $.25(0.6)=0.15$ and net result $14+0-6=8$. The pairs $(0,6)$ and $(6,0)$ both break even; group their masses. Subtract the cost only once and sort signed results numerically.}
\begin{center}
\begin{tabular}{ccc}\toprule
Complete outcome & Probability & $\cTwoValue{X}$ (\cTwoValue{USD}) \\ \midrule
\cTwoOutcome{$(14,0)$} & $\cTwoMassHL{0.15}$ & $\cTwoValue{8}$ \\
\cTwoOutcome{$(14,6)$} & $\cTwoMassHL{0.1}$ & $\cTwoValue{14}$ \\
\cTwoOutcome{$(0,0)$} & $\cTwoMassHL{0.21}$ & $\cTwoValue{-6}$ \\
\cTwoOutcome{$(0,6)$} & $\cTwoMassHL{0.14}$ & $\cTwoValue{0}$ \\
\cTwoOutcome{$(6,0)$} & $\cTwoMassHL{0.24}$ & $\cTwoValue{0}$ \\
\cTwoOutcome{$(6,6)$} & $\cTwoMassHL{0.16}$ & $\cTwoValue{6}$ \\
\bottomrule\end{tabular}
\end{center}
$\cTwoValue{S_X}=\cTwoValue{\{-6,0,6,8,14\}}$; $p_X(x)=0$ elsewhere. The grouped PMF is:
\par\noindent\begin{minipage}[t]{.37\linewidth}\vspace{0pt}\centering
\begin{tabular}{cc}\toprule
$\cTwoValue{x}$ (\cTwoValue{USD}) & $\cTwoMassHL{p_X(x)}$ \\ \midrule
$\cTwoValue{-6}$ & $\cTwoMassHL{0.21}$ \\
$\cTwoValue{0}$ & $\cTwoMassHL{0.38}$ \\
$\cTwoValue{6}$ & $\cTwoMassHL{0.16}$ \\
$\cTwoValue{8}$ & $\cTwoMassHL{0.15}$ \\
$\cTwoValue{14}$ & $\cTwoMassHL{0.1}$ \\
\bottomrule\end{tabular}
\end{minipage}\hfill\begin{minipage}[t]{.60\linewidth}\vspace{0pt}\centering
\GTenPMFGraph[color=cTwoMassColor!20,xhighlight=GTenHLValue,yhighlight=cTwoMassMarker,width=\linewidth,height=4.15cm,xlabel={$\cTwoValue{x}$ (\cTwoValue{USD})},ylabel={$\cTwoMassHL{p_X(x)}$},xstep=1]{-6,0,6,8,14}{0.21,0.38,0.16,0.15,0.1}
\end{minipage}\par
\textbf{PMF check.} All masses are positive and $0.21+0.38+0.16+0.15+0.1=1$. The graph uses the actual numerical spacing; gaps carry no mass.
\task{2}{Accumulate in increasing support order: $0.21;\ 0.59;\ 0.75;\ 0.9;\ 1$. Each new height is the previous height plus the next PMF mass.}
\par\noindent\begin{minipage}[t]{.37\linewidth}\vspace{0pt}\centering
\begin{tabular}{cc}\toprule
Cutoff interval & $\cTwoCDFHL{F_X(x)}$ \\ \midrule
$x<-6$ & $\cTwoCDFHL{0}$ \\
$-6\leq x<0$ & $\cTwoCDFHL{0.21}$ \\
$0\leq x<6$ & $\cTwoCDFHL{0.59}$ \\
$6\leq x<8$ & $\cTwoCDFHL{0.75}$ \\
$8\leq x<14$ & $\cTwoCDFHL{0.9}$ \\
$x\geq 14$ & $\cTwoCDFHL{1}$ \\
\bottomrule\end{tabular}
\end{minipage}\hfill\begin{minipage}[t]{.60\linewidth}\vspace{0pt}\centering
\begin{tikzpicture}\begin{axis}[classCDF,xmin=-8,xmax=16,xtick={-6,0,6,8,14}]
\addplot[cTwoCDFColor,thick,<-] coordinates {(-8,0)(-6,0)};
\addplot[cTwoCDFColor,thick] coordinates {(-6,0.21)(0,0.21)};
\addplot[cTwoCDFColor,thick] coordinates {(0,0.59)(6,0.59)};
\addplot[cTwoCDFColor,thick] coordinates {(6,0.75)(8,0.75)};
\addplot[cTwoCDFColor,thick] coordinates {(8,0.9)(14,0.9)};
\addplot[cTwoCDFColor,thick,->] coordinates {(14,1)(16,1)};
\addplot[only marks,mark=*,cTwoCDFColor,mark size=1.8pt,mark options={fill=white}] coordinates {(-6,0)(0,0.21)(6,0.59)(8,0.75)(14,0.9)};
\addplot[only marks,mark=*,cTwoCDFColor,mark size=1.8pt] coordinates {(-6,0.21)(0,0.59)(6,0.75)(8,0.9)(14,1)};
\end{axis}\end{tikzpicture}
\end{minipage}\par
\textbf{CDF check.} All heights lie in $[0,1]$ and never decrease. Filled left/open right endpoints make the graph right-continuous. Horizontal tails give limits 0 and 1; the CDF reaches 1 at 14 USD. The jumps, $0.21,\ 0.38,\ 0.16,\ 0.15,\ 0.1$, reproduce the PMF.
\textbf{Interpretation.} $F_X(0)=0.59$: the probability that the net result is nonpositive (a loss or break-even) is 0.59.
\vfill
\clearpage
\stageheading{Notation reference}
\noindent\begin{tabularx}{\linewidth}{@{}>{\raggedright\arraybackslash}p{.24\linewidth}>{\raggedright\arraybackslash}p{.16\linewidth}>{\raggedright\arraybackslash}p{.31\linewidth}>{\raggedright\arraybackslash}X@{}}\toprule
Symbol or notation & Concept & Definition & Meaning in the introductory examples\\\midrule
$\cTwoValue{X}$, $\cTwoValue{x}$, $\cTwoValue{S_X}$ & \cTwoValue{Quantity} & $X$ is the recorded numerical quantity; $x$ is a possible value or a real cutoff. State units. $S_X$ contains the distinct values with positive probability. & \textbullet\ Example 1: processing fee (USD).\par\textbullet\ Example 2: late fee (USD).\\
$\cTwoMassHL{p_X(x)}=P(X=x)$ & \cTwoMassHL{PMF} & Add probabilities of mutually exclusive complete outcomes giving $x$. Outside $S_X$, $p_X(x)=0$. Check $p_X(x)\geq0$ and $\sum_{x\in S_X}p_X(x)=1$. & \textbullet\ Example 1: $\cTwoMassHL{p_X(2)}=\cTwoMassHL{0.6}$.\\
$\cTwoCDFHL{F_X(x)}=P(X\leq x)$ & \cTwoCDFHL{CDF} & Sort support, then accumulate: $F_X(x)=\sum_{t\in S_X,\,t\leq x}p_X(t)$. This is defined for every real cutoff. & \textbullet\ Example 1: $\cTwoCDFHL{F_X(2)}=\cTwoCDFHL{0.75}$.\\
$\cTwoCDFHL{F_X(a^-)}$ & \cTwoCDFHL{Left-limit height} & Height immediately before $a$; $F_X(a)-F_X(a^-)=p_X(a)$. The jump adds the mass at $a$. & \textbullet\ Example 1: before $\cTwoValue{2\text{ USD}}$, height $\cTwoCDFHL{0.15}$.\\
\bottomrule\end{tabularx}
\clearpage
\stageheading{Guided Example}
\problemtitle{7}{Community market payments}{A}
\textbf{Payment fees.} A community market stall receives one payment. The only methods are cash, debit card, credit card and mobile wallet, with probabilities .20, .15, .40, .25, respectively. Cash has no fee; either card costs 3 USD; a mobile wallet costs 6 USD. There are no other charges. Record the processing fee.
\task{1}{Define $X$, its units and support; write the PMF table and represent it graphically. Show the outcome mapping and probability calculations, then verify the PMF.}
\textbf{Starter.} Start the mapping with cash $\mapsto0$. Which two methods share a fee?
\worklines{4}
\par\noindent\begin{minipage}[t]{.37\linewidth}\vspace{0pt}\centering
\blanktable{$x$ (USD) & $p_X(x)$}\end{minipage}\hfill\begin{minipage}[t]{.60\linewidth}\vspace{0pt}\centering
\blankgraph{$p_X(x)$}
\end{minipage}\par
\textbf{PMF checks:} \hrulefill\par
\task{2}{Write the CDF table for every real cutoff and represent it graphically. Verify its defining properties and that its jumps reproduce the PMF. Interpret $F_X(3)$ in this context.}
\textbf{Prompt.} Start below the minimum at height 0; add one grouped mass at a time. Include both tails and distinguish open and filled endpoints.
\par\noindent\begin{minipage}[t]{.37\linewidth}\vspace{0pt}\centering
\blanktable{Cutoff interval & $F_X(x)$}\end{minipage}\hfill\begin{minipage}[t]{.60\linewidth}\vspace{0pt}\centering
\blankgraph{$F_X(x)$}
\end{minipage}\par
\textbf{CDF checks and contextual interpretation:}\worklines{3}
\vfill
\clearpage
\problemtitle{8}{Kayak returns}{B}
\textbf{Late fees.} A rental service records the return of one kayak. The exhaustive categories are on time or early, 1--5 whole minutes late, 6--15 whole minutes late, and more than 15 minutes late, with probabilities .25, .20, .30, .25, respectively. The late fee is 0 USD, 4 USD, 4 USD and 8 USD, respectively. No other charges apply. Record the late fee.
\task{1}{Define $X$, its units and support; write the PMF table and represent it graphically. Show the outcome mapping and probability calculations, then verify the PMF.}
\textbf{Starter.} Start with on time $\mapsto0$. Apply the fee schedule, rather than using minutes as $X$.
\worklines{4}
\par\noindent\begin{minipage}[t]{.37\linewidth}\vspace{0pt}\centering
\blanktable{$x$ (USD) & $p_X(x)$}\end{minipage}\hfill\begin{minipage}[t]{.60\linewidth}\vspace{0pt}\centering
\blankgraph{$p_X(x)$}
\end{minipage}\par
\textbf{PMF checks:} \hrulefill\par
\task{2}{Write the CDF table for every real cutoff and represent it graphically. Verify its defining properties and that its jumps reproduce the PMF. Interpret $F_X(4)$ in this context.}
\textbf{Prompt.} Start below the minimum at height 0; add one grouped mass at a time. Include both tails and distinguish open and filled endpoints.
\par\noindent\begin{minipage}[t]{.37\linewidth}\vspace{0pt}\centering
\blanktable{Cutoff interval & $F_X(x)$}\end{minipage}\hfill\begin{minipage}[t]{.60\linewidth}\vspace{0pt}\centering
\blankgraph{$F_X(x)$}
\end{minipage}\par
\textbf{CDF checks and contextual interpretation:}\worklines{3}
\vfill
\clearpage
\problemtitle{9}{Packaging buyback}{C}
\textbf{Uneven buyback payments.} A packaging depot buys one standard-sized bag containing only one material type. The exhaustive types are rejected material, paper, cardboard, plastic, steel and aluminium, with probabilities .12, .18, .15, .20, .10, .25, respectively. Payments are 0 USD for rejected material, 2 USD for either paper or cardboard, 6 USD for either plastic or steel, and 11 USD for aluminium. There are no other payments. Record the buyback payment.
\task{1}{Define $X$, its units and support; write the PMF table and represent it graphically. Show the outcome mapping and probability calculations, then verify the PMF.}
\textbf{Starter.} Start with rejected material $\mapsto0$. Identify both pairs that share a payment.
\worklines{4}
\par\noindent\begin{minipage}[t]{.37\linewidth}\vspace{0pt}\centering
\blanktable{$x$ (USD) & $p_X(x)$}\end{minipage}\hfill\begin{minipage}[t]{.60\linewidth}\vspace{0pt}\centering
\blankgraph{$p_X(x)$}
\end{minipage}\par
\textbf{PMF checks:} \hrulefill\par
\task{2}{Write the CDF table for every real cutoff and represent it graphically. Verify its defining properties and that its jumps reproduce the PMF. Interpret $F_X(6)$ in this context.}
\textbf{Prompt.} Start below the minimum at height 0; add one grouped mass at a time. Include both tails and distinguish open and filled endpoints.
\par\noindent\begin{minipage}[t]{.37\linewidth}\vspace{0pt}\centering
\blanktable{Cutoff interval & $F_X(x)$}\end{minipage}\hfill\begin{minipage}[t]{.60\linewidth}\vspace{0pt}\centering
\blankgraph{$F_X(x)$}
\end{minipage}\par
\textbf{CDF checks and contextual interpretation:}\worklines{3}
\vfill
\clearpage
\problemtitle{10}{Device repair claims}{D}
\textbf{Missing probability and capped payout.} A device repair fund can receive only 0, 1, 2, 3 or 4 claims in one day. Their probabilities, in that order, are .30, .15, a missing probability, .20, .15. It pays 4 USD per claim, capped at 12 USD in total for the day; it pays the smaller of 4 times the count and 12 USD. There are no other payments. Record the daily payout.
\task{1}{Define $X$, its units and support; write the PMF table and represent it graphically. Show the outcome mapping and probability calculations, then verify the PMF.}
\textbf{Starter.} First recover $P(N=2)=1-(.30+.15+.20+.15)$. Then apply the cap.
\worklines{4}
\par\noindent\begin{minipage}[t]{.37\linewidth}\vspace{0pt}\centering
\blanktable{$x$ (USD) & $p_X(x)$}\end{minipage}\hfill\begin{minipage}[t]{.60\linewidth}\vspace{0pt}\centering
\blankgraph{$p_X(x)$}
\end{minipage}\par
\textbf{PMF checks:} \hrulefill\par
\task{2}{Write the CDF table for every real cutoff and represent it graphically. Verify its defining properties and that its jumps reproduce the PMF. Interpret $F_X(8)$ in this context.}
\textbf{Prompt.} Start below the minimum at height 0; add one grouped mass at a time. Include both tails and distinguish open and filled endpoints.
\par\noindent\begin{minipage}[t]{.37\linewidth}\vspace{0pt}\centering
\blanktable{Cutoff interval & $F_X(x)$}\end{minipage}\hfill\begin{minipage}[t]{.60\linewidth}\vspace{0pt}\centering
\blankgraph{$F_X(x)$}
\end{minipage}\par
\textbf{CDF checks and contextual interpretation:}\worklines{3}
\vfill
\clearpage
\problemtitle{11}{Mug sales}{E}
\textbf{Independent purchases and discount.} A school enterprise offers one mug to each of three visitors. Independently, each buys one with probability 0.6 or passes with probability 0.4, regardless of the discount. Each item costs 6 USD. Only if all three buy, subtract a 4 USD discount from the combined bill. Otherwise there is no discount. There are no other sales or charges. Record total revenue.
\task{1}{Define $X$, its units and support; write the PMF table and represent it graphically. Show the outcome mapping and probability calculations, then verify the PMF.}
\textbf{Starter.} Write B/N in visitor order. Start with $P(\mathrm{NNN})=(0.4)^3$. List all eight sequences and apply the discount only to BBB.
\par\noindent\begin{tabularx}{\linewidth}{|Y|Y|Y|Y|Y|Y|}\hline Outcome & Probability & $X$ (USD) & Outcome & Probability & $X$ (USD)\\\hline
\rule{0pt}{.63cm}& & & & & \\\hline
\rule{0pt}{.63cm}& & & & & \\\hline
\rule{0pt}{.63cm}& & & & & \\\hline
\rule{0pt}{.63cm}& & & & & \\\hline
\end{tabularx}
\par\noindent\begin{minipage}[t]{.37\linewidth}\vspace{0pt}\centering
\blanktable{$x$ (USD) & $p_X(x)$}\end{minipage}\hfill\begin{minipage}[t]{.60\linewidth}\vspace{0pt}\centering
\blankgraph{$p_X(x)$}
\end{minipage}\par
\textbf{PMF checks:} \hrulefill\par
\task{2}{Write the CDF table for every real cutoff and represent it graphically. Verify its defining properties and that its jumps reproduce the PMF. Interpret $F_X(12)$ in this context.}
\textbf{Prompt.} Start below the minimum at height 0; add one grouped mass at a time. Include both tails and distinguish open and filled endpoints.
\par\noindent\begin{minipage}[t]{.37\linewidth}\vspace{0pt}\centering
\blanktable{Cutoff interval & $F_X(x)$}\end{minipage}\hfill\begin{minipage}[t]{.60\linewidth}\vspace{0pt}\centering
\blankgraph{$F_X(x)$}
\end{minipage}\par
\textbf{CDF checks and contextual interpretation:}\worklines{3}
\vfill
\clearpage
\problemtitle{12}{Two-stage garden fund}{F}
\textbf{Signed two-stage net results.} A hypothetical classroom fund receives two cash returns. The first is 12 USD, 0 USD or 5 USD with probabilities .40, .25, .35, respectively. Independently, the second is 0 USD with probability 0.7 or 5 USD with probability 0.3. These are all possible returns. A fixed participation cost of 5 USD is charged once. There are no other receipts or charges. Record the sum of both returns minus the cost; a negative net result means a loss.
\task{1}{Define $X$, its units and support; write the PMF table and represent it graphically. Show the outcome mapping and probability calculations, then verify the PMF.}
\textbf{Starter.} Start with $(R_1,R_2)=(12,0)$: probability $.40(0.7)$, net $12+0-5$. List all six ordered pairs.
\par\noindent\begin{tabularx}{\linewidth}{|Y|Y|Y|Y|Y|Y|}\hline Outcome & Probability & $X$ (USD) & Outcome & Probability & $X$ (USD)\\\hline
\rule{0pt}{.63cm}& & & & & \\\hline
\rule{0pt}{.63cm}& & & & & \\\hline
\rule{0pt}{.63cm}& & & & & \\\hline
\end{tabularx}
\par\noindent\begin{minipage}[t]{.37\linewidth}\vspace{0pt}\centering
\blanktable{$x$ (USD) & $p_X(x)$}\end{minipage}\hfill\begin{minipage}[t]{.60\linewidth}\vspace{0pt}\centering
\blankgraph{$p_X(x)$}
\end{minipage}\par
\textbf{PMF checks:} \hrulefill\par
\task{2}{Write the CDF table for every real cutoff and represent it graphically. Verify its defining properties and that its jumps reproduce the PMF. Interpret $F_X(0)$ in this context.}
\textbf{Prompt.} Start below the minimum at height 0; add one grouped mass at a time. Include both tails and distinguish open and filled endpoints.
\par\noindent\begin{minipage}[t]{.37\linewidth}\vspace{0pt}\centering
\blanktable{Cutoff interval & $F_X(x)$}\end{minipage}\hfill\begin{minipage}[t]{.60\linewidth}\vspace{0pt}\centering
\blankgraph{$F_X(x)$}
\end{minipage}\par
\textbf{CDF checks and contextual interpretation:}\worklines{3}
\vfill
\clearpage
\stageheading{Independent Practice}
\problemtitle{13}{Theatre ticket payments}{A}
\textbf{Payment fees.} A theatre ticket booth receives one payment. The only methods are cash, debit card, credit card and mobile wallet, with probabilities .30, .20, .15, .35, respectively. Cash has no fee; either card costs 4 USD; a mobile wallet costs 8 USD. There are no other charges. Record the processing fee.
\task{1}{Define $X$, its units and support; write the PMF table and represent it graphically. Show the outcome mapping and probability calculations, then verify the PMF.}
\textbf{Calculation area:} include every complete outcome and its probability before grouping.
\worklines{4}
\par\noindent\begin{minipage}[t]{.37\linewidth}\vspace{0pt}\centering
\blanktable{$x$ (USD) & $p_X(x)$}\end{minipage}\hfill\begin{minipage}[t]{.60\linewidth}\vspace{0pt}\centering
\blankgraph{$p_X(x)$}
\end{minipage}\par
\textbf{PMF checks:} \hrulefill\par
\task{2}{Write the CDF table for every real cutoff and represent it graphically. Verify its defining properties and that its jumps reproduce the PMF. Interpret $F_X(4)$ in this context.}
\par\noindent\begin{minipage}[t]{.37\linewidth}\vspace{0pt}\centering
\blanktable{Cutoff interval & $F_X(x)$}\end{minipage}\hfill\begin{minipage}[t]{.60\linewidth}\vspace{0pt}\centering
\blankgraph{$F_X(x)$}
\end{minipage}\par
\textbf{CDF checks and contextual interpretation:}\worklines{3}
\vfill
\clearpage
\problemtitle{14}{Sports stall payments}{A}
\textbf{Payment fees.} A sports stall receives one payment. The only methods are cash, debit card, credit card and mobile wallet, with probabilities .12, .28, .22, .38, respectively. Cash has no fee; either card costs 5 USD; a mobile wallet costs 10 USD. There are no other charges. Record the processing fee.
\task{1}{Define $X$, its units and support; write the PMF table and represent it graphically. Show the outcome mapping and probability calculations, then verify the PMF.}
\textbf{Calculation area:} include every complete outcome and its probability before grouping.
\worklines{4}
\par\noindent\begin{minipage}[t]{.37\linewidth}\vspace{0pt}\centering
\blanktable{$x$ (USD) & $p_X(x)$}\end{minipage}\hfill\begin{minipage}[t]{.60\linewidth}\vspace{0pt}\centering
\blankgraph{$p_X(x)$}
\end{minipage}\par
\textbf{PMF checks:} \hrulefill\par
\task{2}{Write the CDF table for every real cutoff and represent it graphically. Verify its defining properties and that its jumps reproduce the PMF. Interpret $F_X(5)$ in this context.}
\par\noindent\begin{minipage}[t]{.37\linewidth}\vspace{0pt}\centering
\blanktable{Cutoff interval & $F_X(x)$}\end{minipage}\hfill\begin{minipage}[t]{.60\linewidth}\vspace{0pt}\centering
\blankgraph{$F_X(x)$}
\end{minipage}\par
\textbf{CDF checks and contextual interpretation:}\worklines{3}
\vfill
\clearpage
\problemtitle{15}{Camera returns}{B}
\textbf{Late fees.} A rental service records the return of one camera. The exhaustive categories are on time or early, 1--5 whole minutes late, 6--15 whole minutes late, and more than 15 minutes late, with probabilities .18, .22, .35, .25, respectively. The late fee is 0 USD, 3 USD, 3 USD and 6 USD, respectively. No other charges apply. Record the late fee.
\task{1}{Define $X$, its units and support; write the PMF table and represent it graphically. Show the outcome mapping and probability calculations, then verify the PMF.}
\textbf{Calculation area:} include every complete outcome and its probability before grouping.
\worklines{4}
\par\noindent\begin{minipage}[t]{.37\linewidth}\vspace{0pt}\centering
\blanktable{$x$ (USD) & $p_X(x)$}\end{minipage}\hfill\begin{minipage}[t]{.60\linewidth}\vspace{0pt}\centering
\blankgraph{$p_X(x)$}
\end{minipage}\par
\textbf{PMF checks:} \hrulefill\par
\task{2}{Write the CDF table for every real cutoff and represent it graphically. Verify its defining properties and that its jumps reproduce the PMF. Interpret $F_X(3)$ in this context.}
\par\noindent\begin{minipage}[t]{.37\linewidth}\vspace{0pt}\centering
\blanktable{Cutoff interval & $F_X(x)$}\end{minipage}\hfill\begin{minipage}[t]{.60\linewidth}\vspace{0pt}\centering
\blankgraph{$F_X(x)$}
\end{minipage}\par
\textbf{CDF checks and contextual interpretation:}\worklines{3}
\vfill
\clearpage
\problemtitle{16}{Tool returns}{B}
\textbf{Late fees.} A rental service records the return of one tool. The exhaustive categories are on time or early, 1--5 whole minutes late, 6--15 whole minutes late, and more than 15 minutes late, with probabilities .35, .15, .20, .30, respectively. The late fee is 0 USD, 5 USD, 5 USD and 10 USD, respectively. No other charges apply. Record the late fee.
\task{1}{Define $X$, its units and support; write the PMF table and represent it graphically. Show the outcome mapping and probability calculations, then verify the PMF.}
\textbf{Calculation area:} include every complete outcome and its probability before grouping.
\worklines{4}
\par\noindent\begin{minipage}[t]{.37\linewidth}\vspace{0pt}\centering
\blanktable{$x$ (USD) & $p_X(x)$}\end{minipage}\hfill\begin{minipage}[t]{.60\linewidth}\vspace{0pt}\centering
\blankgraph{$p_X(x)$}
\end{minipage}\par
\textbf{PMF checks:} \hrulefill\par
\task{2}{Write the CDF table for every real cutoff and represent it graphically. Verify its defining properties and that its jumps reproduce the PMF. Interpret $F_X(5)$ in this context.}
\par\noindent\begin{minipage}[t]{.37\linewidth}\vspace{0pt}\centering
\blanktable{Cutoff interval & $F_X(x)$}\end{minipage}\hfill\begin{minipage}[t]{.60\linewidth}\vspace{0pt}\centering
\blankgraph{$F_X(x)$}
\end{minipage}\par
\textbf{CDF checks and contextual interpretation:}\worklines{3}
\vfill
\clearpage
\problemtitle{17}{Workshop scrap buyback}{C}
\textbf{Uneven buyback payments.} A workshop scrap centre buys one standard-sized bag containing only one material type. The exhaustive types are rejected material, paper, cardboard, plastic, steel and aluminium, with probabilities .20, .12, .18, .15, .10, .25, respectively. Payments are 0 USD for rejected material, 4 USD for either paper or cardboard, 9 USD for either plastic or steel, and 15 USD for aluminium. There are no other payments. Record the buyback payment.
\task{1}{Define $X$, its units and support; write the PMF table and represent it graphically. Show the outcome mapping and probability calculations, then verify the PMF.}
\textbf{Calculation area:} include every complete outcome and its probability before grouping.
\worklines{4}
\par\noindent\begin{minipage}[t]{.37\linewidth}\vspace{0pt}\centering
\blanktable{$x$ (USD) & $p_X(x)$}\end{minipage}\hfill\begin{minipage}[t]{.60\linewidth}\vspace{0pt}\centering
\blankgraph{$p_X(x)$}
\end{minipage}\par
\textbf{PMF checks:} \hrulefill\par
\task{2}{Write the CDF table for every real cutoff and represent it graphically. Verify its defining properties and that its jumps reproduce the PMF. Interpret $F_X(9)$ in this context.}
\par\noindent\begin{minipage}[t]{.37\linewidth}\vspace{0pt}\centering
\blanktable{Cutoff interval & $F_X(x)$}\end{minipage}\hfill\begin{minipage}[t]{.60\linewidth}\vspace{0pt}\centering
\blankgraph{$F_X(x)$}
\end{minipage}\par
\textbf{CDF checks and contextual interpretation:}\worklines{3}
\vfill
\clearpage
\problemtitle{18}{Community salvage buyback}{C}
\textbf{Uneven buyback payments.} A community salvage stall buys one standard-sized bag containing only one material type. The exhaustive types are rejected material, paper, cardboard, plastic, steel and aluminium, with probabilities .08, .17, .15, .25, .15, .20, respectively. Payments are 0 USD for rejected material, 3 USD for either paper or cardboard, 8 USD for either plastic or steel, and 14 USD for aluminium. There are no other payments. Record the buyback payment.
\task{1}{Define $X$, its units and support; write the PMF table and represent it graphically. Show the outcome mapping and probability calculations, then verify the PMF.}
\textbf{Calculation area:} include every complete outcome and its probability before grouping.
\worklines{4}
\par\noindent\begin{minipage}[t]{.37\linewidth}\vspace{0pt}\centering
\blanktable{$x$ (USD) & $p_X(x)$}\end{minipage}\hfill\begin{minipage}[t]{.60\linewidth}\vspace{0pt}\centering
\blankgraph{$p_X(x)$}
\end{minipage}\par
\textbf{PMF checks:} \hrulefill\par
\task{2}{Write the CDF table for every real cutoff and represent it graphically. Verify its defining properties and that its jumps reproduce the PMF. Interpret $F_X(8)$ in this context.}
\par\noindent\begin{minipage}[t]{.37\linewidth}\vspace{0pt}\centering
\blanktable{Cutoff interval & $F_X(x)$}\end{minipage}\hfill\begin{minipage}[t]{.60\linewidth}\vspace{0pt}\centering
\blankgraph{$F_X(x)$}
\end{minipage}\par
\textbf{CDF checks and contextual interpretation:}\worklines{3}
\vfill
\clearpage
\problemtitle{19}{Delivery loss claims}{D}
\textbf{Missing probability and capped payout.} A delivery insurer can receive only 0, 1, 2, 3 or 4 claims in one day. Their probabilities, in that order, are .15, .20, a missing probability, .20, .10. It pays 5 USD per claim, capped at 15 USD in total for the day; it pays the smaller of 5 times the count and 15 USD. There are no other payments. Record the daily payout.
\task{1}{Define $X$, its units and support; write the PMF table and represent it graphically. Show the outcome mapping and probability calculations, then verify the PMF.}
\textbf{Calculation area:} include every complete outcome and its probability before grouping.
\worklines{4}
\par\noindent\begin{minipage}[t]{.37\linewidth}\vspace{0pt}\centering
\blanktable{$x$ (USD) & $p_X(x)$}\end{minipage}\hfill\begin{minipage}[t]{.60\linewidth}\vspace{0pt}\centering
\blankgraph{$p_X(x)$}
\end{minipage}\par
\textbf{PMF checks:} \hrulefill\par
\task{2}{Write the CDF table for every real cutoff and represent it graphically. Verify its defining properties and that its jumps reproduce the PMF. Interpret $F_X(10)$ in this context.}
\par\noindent\begin{minipage}[t]{.37\linewidth}\vspace{0pt}\centering
\blanktable{Cutoff interval & $F_X(x)$}\end{minipage}\hfill\begin{minipage}[t]{.60\linewidth}\vspace{0pt}\centering
\blankgraph{$F_X(x)$}
\end{minipage}\par
\textbf{CDF checks and contextual interpretation:}\worklines{3}
\vfill
\clearpage
\problemtitle{20}{Event cancellation claims}{D}
\textbf{Missing probability and capped payout.} An event protection fund can receive only 0, 1, 2, 3 or 4 claims in one day. Their probabilities, in that order, are .22, .18, a missing probability, .17, .15. It pays 6 USD per claim, capped at 18 USD in total for the day; it pays the smaller of 6 times the count and 18 USD. There are no other payments. Record the daily payout.
\task{1}{Define $X$, its units and support; write the PMF table and represent it graphically. Show the outcome mapping and probability calculations, then verify the PMF.}
\textbf{Calculation area:} include every complete outcome and its probability before grouping.
\worklines{4}
\par\noindent\begin{minipage}[t]{.37\linewidth}\vspace{0pt}\centering
\blanktable{$x$ (USD) & $p_X(x)$}\end{minipage}\hfill\begin{minipage}[t]{.60\linewidth}\vspace{0pt}\centering
\blankgraph{$p_X(x)$}
\end{minipage}\par
\textbf{PMF checks:} \hrulefill\par
\task{2}{Write the CDF table for every real cutoff and represent it graphically. Verify its defining properties and that its jumps reproduce the PMF. Interpret $F_X(12)$ in this context.}
\par\noindent\begin{minipage}[t]{.37\linewidth}\vspace{0pt}\centering
\blanktable{Cutoff interval & $F_X(x)$}\end{minipage}\hfill\begin{minipage}[t]{.60\linewidth}\vspace{0pt}\centering
\blankgraph{$F_X(x)$}
\end{minipage}\par
\textbf{CDF checks and contextual interpretation:}\worklines{3}
\vfill
\clearpage
\problemtitle{21}{Tote bag sales}{E}
\textbf{Independent purchases and discount.} A school enterprise offers one tote bag to each of three visitors. Independently, each buys one with probability 0.3 or passes with probability 0.7, regardless of the discount. Each item costs 7 USD. Only if all three buy, subtract a 5 USD discount from the combined bill. Otherwise there is no discount. There are no other sales or charges. Record total revenue.
\task{1}{Define $X$, its units and support; write the PMF table and represent it graphically. Show the outcome mapping and probability calculations, then verify the PMF.}
\textbf{Calculation area:} include every complete outcome and its probability before grouping.
\par\noindent\begin{tabularx}{\linewidth}{|Y|Y|Y|Y|Y|Y|}\hline Outcome & Probability & $X$ (USD) & Outcome & Probability & $X$ (USD)\\\hline
\rule{0pt}{.63cm}& & & & & \\\hline
\rule{0pt}{.63cm}& & & & & \\\hline
\rule{0pt}{.63cm}& & & & & \\\hline
\rule{0pt}{.63cm}& & & & & \\\hline
\end{tabularx}
\par\noindent\begin{minipage}[t]{.37\linewidth}\vspace{0pt}\centering
\blanktable{$x$ (USD) & $p_X(x)$}\end{minipage}\hfill\begin{minipage}[t]{.60\linewidth}\vspace{0pt}\centering
\blankgraph{$p_X(x)$}
\end{minipage}\par
\textbf{PMF checks:} \hrulefill\par
\task{2}{Write the CDF table for every real cutoff and represent it graphically. Verify its defining properties and that its jumps reproduce the PMF. Interpret $F_X(14)$ in this context.}
\par\noindent\begin{minipage}[t]{.37\linewidth}\vspace{0pt}\centering
\blanktable{Cutoff interval & $F_X(x)$}\end{minipage}\hfill\begin{minipage}[t]{.60\linewidth}\vspace{0pt}\centering
\blankgraph{$F_X(x)$}
\end{minipage}\par
\textbf{CDF checks and contextual interpretation:}\worklines{3}
\vfill
\clearpage
\problemtitle{22}{Poster sales}{E}
\textbf{Independent purchases and discount.} A school enterprise offers one poster to each of three visitors. Independently, each buys one with probability 0.7 or passes with probability 0.3, regardless of the discount. Each item costs 8 USD. Only if all three buy, subtract a 3 USD discount from the combined bill. Otherwise there is no discount. There are no other sales or charges. Record total revenue.
\task{1}{Define $X$, its units and support; write the PMF table and represent it graphically. Show the outcome mapping and probability calculations, then verify the PMF.}
\textbf{Calculation area:} include every complete outcome and its probability before grouping.
\par\noindent\begin{tabularx}{\linewidth}{|Y|Y|Y|Y|Y|Y|}\hline Outcome & Probability & $X$ (USD) & Outcome & Probability & $X$ (USD)\\\hline
\rule{0pt}{.63cm}& & & & & \\\hline
\rule{0pt}{.63cm}& & & & & \\\hline
\rule{0pt}{.63cm}& & & & & \\\hline
\rule{0pt}{.63cm}& & & & & \\\hline
\end{tabularx}
\par\noindent\begin{minipage}[t]{.37\linewidth}\vspace{0pt}\centering
\blanktable{$x$ (USD) & $p_X(x)$}\end{minipage}\hfill\begin{minipage}[t]{.60\linewidth}\vspace{0pt}\centering
\blankgraph{$p_X(x)$}
\end{minipage}\par
\textbf{PMF checks:} \hrulefill\par
\task{2}{Write the CDF table for every real cutoff and represent it graphically. Verify its defining properties and that its jumps reproduce the PMF. Interpret $F_X(16)$ in this context.}
\par\noindent\begin{minipage}[t]{.37\linewidth}\vspace{0pt}\centering
\blanktable{Cutoff interval & $F_X(x)$}\end{minipage}\hfill\begin{minipage}[t]{.60\linewidth}\vspace{0pt}\centering
\blankgraph{$F_X(x)$}
\end{minipage}\par
\textbf{CDF checks and contextual interpretation:}\worklines{3}
\vfill
\clearpage
\problemtitle{23}{Two-stage media fund}{F}
\textbf{Signed two-stage net results.} A hypothetical classroom fund receives two cash returns. The first is 17 USD, 0 USD or 7 USD with probabilities .20, .30, .50, respectively. Independently, the second is 0 USD with probability 0.4 or 7 USD with probability 0.6. These are all possible returns. A fixed participation cost of 7 USD is charged once. There are no other receipts or charges. Record the sum of both returns minus the cost; a negative net result means a loss.
\task{1}{Define $X$, its units and support; write the PMF table and represent it graphically. Show the outcome mapping and probability calculations, then verify the PMF.}
\textbf{Calculation area:} include every complete outcome and its probability before grouping.
\par\noindent\begin{tabularx}{\linewidth}{|Y|Y|Y|Y|Y|Y|}\hline Outcome & Probability & $X$ (USD) & Outcome & Probability & $X$ (USD)\\\hline
\rule{0pt}{.63cm}& & & & & \\\hline
\rule{0pt}{.63cm}& & & & & \\\hline
\rule{0pt}{.63cm}& & & & & \\\hline
\end{tabularx}
\par\noindent\begin{minipage}[t]{.37\linewidth}\vspace{0pt}\centering
\blanktable{$x$ (USD) & $p_X(x)$}\end{minipage}\hfill\begin{minipage}[t]{.60\linewidth}\vspace{0pt}\centering
\blankgraph{$p_X(x)$}
\end{minipage}\par
\textbf{PMF checks:} \hrulefill\par
\task{2}{Write the CDF table for every real cutoff and represent it graphically. Verify its defining properties and that its jumps reproduce the PMF. Interpret $F_X(0)$ in this context.}
\par\noindent\begin{minipage}[t]{.37\linewidth}\vspace{0pt}\centering
\blanktable{Cutoff interval & $F_X(x)$}\end{minipage}\hfill\begin{minipage}[t]{.60\linewidth}\vspace{0pt}\centering
\blankgraph{$F_X(x)$}
\end{minipage}\par
\textbf{CDF checks and contextual interpretation:}\worklines{3}
\vfill
\clearpage
\problemtitle{24}{Two-stage sports fund}{F}
\textbf{Signed two-stage net results.} A hypothetical classroom fund receives two cash returns. The first is 19 USD, 0 USD or 8 USD with probabilities .35, .15, .50, respectively. Independently, the second is 0 USD with probability 0.8 or 8 USD with probability 0.2. These are all possible returns. A fixed participation cost of 8 USD is charged once. There are no other receipts or charges. Record the sum of both returns minus the cost; a negative net result means a loss.
\task{1}{Define $X$, its units and support; write the PMF table and represent it graphically. Show the outcome mapping and probability calculations, then verify the PMF.}
\textbf{Calculation area:} include every complete outcome and its probability before grouping.
\par\noindent\begin{tabularx}{\linewidth}{|Y|Y|Y|Y|Y|Y|}\hline Outcome & Probability & $X$ (USD) & Outcome & Probability & $X$ (USD)\\\hline
\rule{0pt}{.63cm}& & & & & \\\hline
\rule{0pt}{.63cm}& & & & & \\\hline
\rule{0pt}{.63cm}& & & & & \\\hline
\end{tabularx}
\par\noindent\begin{minipage}[t]{.37\linewidth}\vspace{0pt}\centering
\blanktable{$x$ (USD) & $p_X(x)$}\end{minipage}\hfill\begin{minipage}[t]{.60\linewidth}\vspace{0pt}\centering
\blankgraph{$p_X(x)$}
\end{minipage}\par
\textbf{PMF checks:} \hrulefill\par
\task{2}{Write the CDF table for every real cutoff and represent it graphically. Verify its defining properties and that its jumps reproduce the PMF. Interpret $F_X(0)$ in this context.}
\par\noindent\begin{minipage}[t]{.37\linewidth}\vspace{0pt}\centering
\blanktable{Cutoff interval & $F_X(x)$}\end{minipage}\hfill\begin{minipage}[t]{.60\linewidth}\vspace{0pt}\centering
\blankgraph{$F_X(x)$}
\end{minipage}\par
\textbf{CDF checks and contextual interpretation:}\worklines{3}
\vfill
\ifprintsolutions
\clearpage
{\centering\large\bfseries Solutions\par}
\problemtitle{7}{Community market payments}{A}
\textbf{Payment fees.} A community market stall receives one payment. The only methods are cash, debit card, credit card and mobile wallet, with probabilities .20, .15, .40, .25, respectively. Cash has no fee; either card costs 3 USD; a mobile wallet costs 6 USD. There are no other charges. Record the processing fee.
\task{1}{Let $\cTwoValue{X}$ be the processing fee for one payment, in USD. Cash maps to 0, either card to 3, and wallet to 6. The two card events are mutually exclusive, so their masses add: $p_X(3)=.15+.40=0.55$.}
\begin{center}
\begin{tabular}{ccc}\toprule
Complete outcome & Probability & $\cTwoValue{X}$ (\cTwoValue{USD}) \\ \midrule
\cTwoOutcome{Cash} & $\cTwoMassHL{0.2}$ & $\cTwoValue{0}$ \\
\cTwoOutcome{Debit card} & $\cTwoMassHL{0.15}$ & $\cTwoValue{3}$ \\
\cTwoOutcome{Credit card} & $\cTwoMassHL{0.4}$ & $\cTwoValue{3}$ \\
\cTwoOutcome{Wallet} & $\cTwoMassHL{0.25}$ & $\cTwoValue{6}$ \\
\bottomrule\end{tabular}
\end{center}
$\cTwoValue{S_X}=\cTwoValue{\{0,3,6\}}$; $p_X(x)=0$ elsewhere. The grouped PMF is:
\par\noindent\begin{minipage}[t]{.37\linewidth}\vspace{0pt}\centering
\begin{tabular}{cc}\toprule
$\cTwoValue{x}$ (\cTwoValue{USD}) & $\cTwoMassHL{p_X(x)}$ \\ \midrule
$\cTwoValue{0}$ & $\cTwoMassHL{0.2}$ \\
$\cTwoValue{3}$ & $\cTwoMassHL{0.55}$ \\
$\cTwoValue{6}$ & $\cTwoMassHL{0.25}$ \\
\bottomrule\end{tabular}
\end{minipage}\hfill\begin{minipage}[t]{.60\linewidth}\vspace{0pt}\centering
\GTenPMFGraph[color=cTwoMassColor!20,xhighlight=GTenHLValue,yhighlight=cTwoMassMarker,width=\linewidth,height=4.15cm,xlabel={$\cTwoValue{x}$ (\cTwoValue{USD})},ylabel={$\cTwoMassHL{p_X(x)}$},xstep=1]{0,3,6}{0.2,0.55,0.25}
\end{minipage}\par
\textbf{PMF check.} All masses are positive and $0.2+0.55+0.25=1$. The graph uses the actual numerical spacing; gaps carry no mass.
\task{2}{Accumulate in increasing support order: $0.2;\ 0.75;\ 1$. Each new height is the previous height plus the next PMF mass.}
\par\noindent\begin{minipage}[t]{.37\linewidth}\vspace{0pt}\centering
\begin{tabular}{cc}\toprule
Cutoff interval & $\cTwoCDFHL{F_X(x)}$ \\ \midrule
$x<0$ & $\cTwoCDFHL{0}$ \\
$0\leq x<3$ & $\cTwoCDFHL{0.2}$ \\
$3\leq x<6$ & $\cTwoCDFHL{0.75}$ \\
$x\geq 6$ & $\cTwoCDFHL{1}$ \\
\bottomrule\end{tabular}
\end{minipage}\hfill\begin{minipage}[t]{.60\linewidth}\vspace{0pt}\centering
\begin{tikzpicture}\begin{axis}[classCDF,xmin=-3,xmax=9,xtick={0,3,6}]
\addplot[cTwoCDFColor,thick,<-] coordinates {(-3,0)(0,0)};
\addplot[cTwoCDFColor,thick] coordinates {(0,0.2)(3,0.2)};
\addplot[cTwoCDFColor,thick] coordinates {(3,0.75)(6,0.75)};
\addplot[cTwoCDFColor,thick,->] coordinates {(6,1)(9,1)};
\addplot[only marks,mark=*,cTwoCDFColor,mark size=1.8pt,mark options={fill=white}] coordinates {(0,0)(3,0.2)(6,0.75)};
\addplot[only marks,mark=*,cTwoCDFColor,mark size=1.8pt] coordinates {(0,0.2)(3,0.75)(6,1)};
\end{axis}\end{tikzpicture}
\end{minipage}\par
\textbf{CDF check.} All heights lie in $[0,1]$ and never decrease. Filled left/open right endpoints make the graph right-continuous. Horizontal tails give limits 0 and 1; the CDF reaches 1 at 6 USD. The jumps, $0.2,\ 0.55,\ 0.25$, reproduce the PMF.
\textbf{Interpretation.} $F_X(3)=0.75$: the probability that the fee is at most 3 USD (cash or either card) is 0.75.
\vfill
\clearpage
\problemtitle{8}{Kayak returns}{B}
\textbf{Late fees.} A rental service records the return of one kayak. The exhaustive categories are on time or early, 1--5 whole minutes late, 6--15 whole minutes late, and more than 15 minutes late, with probabilities .25, .20, .30, .25, respectively. The late fee is 0 USD, 4 USD, 4 USD and 8 USD, respectively. No other charges apply. Record the late fee.
\task{1}{Let $\cTwoValue{X}$ be the late fee for one return, in USD. The two shorter delays both map to 4 USD. Their mutually exclusive masses add: $p_X(4)=.20+.30=0.5$. On-time returns map to 0 and the longest delay maps to 8.}
\begin{center}
\begin{tabular}{ccc}\toprule
Complete outcome & Probability & $\cTwoValue{X}$ (\cTwoValue{USD}) \\ \midrule
\cTwoOutcome{On time/early} & $\cTwoMassHL{0.25}$ & $\cTwoValue{0}$ \\
\cTwoOutcome{1--5 min} & $\cTwoMassHL{0.2}$ & $\cTwoValue{4}$ \\
\cTwoOutcome{6--15 min} & $\cTwoMassHL{0.3}$ & $\cTwoValue{4}$ \\
\cTwoOutcome{Over 15 min} & $\cTwoMassHL{0.25}$ & $\cTwoValue{8}$ \\
\bottomrule\end{tabular}
\end{center}
$\cTwoValue{S_X}=\cTwoValue{\{0,4,8\}}$; $p_X(x)=0$ elsewhere. The grouped PMF is:
\par\noindent\begin{minipage}[t]{.37\linewidth}\vspace{0pt}\centering
\begin{tabular}{cc}\toprule
$\cTwoValue{x}$ (\cTwoValue{USD}) & $\cTwoMassHL{p_X(x)}$ \\ \midrule
$\cTwoValue{0}$ & $\cTwoMassHL{0.25}$ \\
$\cTwoValue{4}$ & $\cTwoMassHL{0.5}$ \\
$\cTwoValue{8}$ & $\cTwoMassHL{0.25}$ \\
\bottomrule\end{tabular}
\end{minipage}\hfill\begin{minipage}[t]{.60\linewidth}\vspace{0pt}\centering
\GTenPMFGraph[color=cTwoMassColor!20,xhighlight=GTenHLValue,yhighlight=cTwoMassMarker,width=\linewidth,height=4.15cm,xlabel={$\cTwoValue{x}$ (\cTwoValue{USD})},ylabel={$\cTwoMassHL{p_X(x)}$},xstep=1]{0,4,8}{0.25,0.5,0.25}
\end{minipage}\par
\textbf{PMF check.} All masses are positive and $0.25+0.5+0.25=1$. The graph uses the actual numerical spacing; gaps carry no mass.
\task{2}{Accumulate in increasing support order: $0.25;\ 0.75;\ 1$. Each new height is the previous height plus the next PMF mass.}
\par\noindent\begin{minipage}[t]{.37\linewidth}\vspace{0pt}\centering
\begin{tabular}{cc}\toprule
Cutoff interval & $\cTwoCDFHL{F_X(x)}$ \\ \midrule
$x<0$ & $\cTwoCDFHL{0}$ \\
$0\leq x<4$ & $\cTwoCDFHL{0.25}$ \\
$4\leq x<8$ & $\cTwoCDFHL{0.75}$ \\
$x\geq 8$ & $\cTwoCDFHL{1}$ \\
\bottomrule\end{tabular}
\end{minipage}\hfill\begin{minipage}[t]{.60\linewidth}\vspace{0pt}\centering
\begin{tikzpicture}\begin{axis}[classCDF,xmin=-4,xmax=12,xtick={0,4,8}]
\addplot[cTwoCDFColor,thick,<-] coordinates {(-4,0)(0,0)};
\addplot[cTwoCDFColor,thick] coordinates {(0,0.25)(4,0.25)};
\addplot[cTwoCDFColor,thick] coordinates {(4,0.75)(8,0.75)};
\addplot[cTwoCDFColor,thick,->] coordinates {(8,1)(12,1)};
\addplot[only marks,mark=*,cTwoCDFColor,mark size=1.8pt,mark options={fill=white}] coordinates {(0,0)(4,0.25)(8,0.75)};
\addplot[only marks,mark=*,cTwoCDFColor,mark size=1.8pt] coordinates {(0,0.25)(4,0.75)(8,1)};
\end{axis}\end{tikzpicture}
\end{minipage}\par
\textbf{CDF check.} All heights lie in $[0,1]$ and never decrease. Filled left/open right endpoints make the graph right-continuous. Horizontal tails give limits 0 and 1; the CDF reaches 1 at 8 USD. The jumps, $0.25,\ 0.5,\ 0.25$, reproduce the PMF.
\textbf{Interpretation.} $F_X(4)=0.75$: the probability that the late fee is at most 4 USD (the return is no more than 15 minutes late) is 0.75.
\vfill
\clearpage
\problemtitle{9}{Packaging buyback}{C}
\textbf{Uneven buyback payments.} A packaging depot buys one standard-sized bag containing only one material type. The exhaustive types are rejected material, paper, cardboard, plastic, steel and aluminium, with probabilities .12, .18, .15, .20, .10, .25, respectively. Payments are 0 USD for rejected material, 2 USD for either paper or cardboard, 6 USD for either plastic or steel, and 11 USD for aluminium. There are no other payments. Record the buyback payment.
\task{1}{Let $\cTwoValue{X}$ be the buyback payment for one bag, in USD. Paper and cardboard share 2 USD, so their probabilities add to 0.33. Plastic and steel share 6 USD, giving 0.3. Keep the actual gaps 2, 4 and 5 USD on both graphs.}
\begin{center}
\begin{tabular}{ccc}\toprule
Complete outcome & Probability & $\cTwoValue{X}$ (\cTwoValue{USD}) \\ \midrule
\cTwoOutcome{Rejected} & $\cTwoMassHL{0.12}$ & $\cTwoValue{0}$ \\
\cTwoOutcome{Paper} & $\cTwoMassHL{0.18}$ & $\cTwoValue{2}$ \\
\cTwoOutcome{Cardboard} & $\cTwoMassHL{0.15}$ & $\cTwoValue{2}$ \\
\cTwoOutcome{Plastic} & $\cTwoMassHL{0.2}$ & $\cTwoValue{6}$ \\
\cTwoOutcome{Steel} & $\cTwoMassHL{0.1}$ & $\cTwoValue{6}$ \\
\cTwoOutcome{Aluminium} & $\cTwoMassHL{0.25}$ & $\cTwoValue{11}$ \\
\bottomrule\end{tabular}
\end{center}
$\cTwoValue{S_X}=\cTwoValue{\{0,2,6,11\}}$; $p_X(x)=0$ elsewhere. The grouped PMF is:
\par\noindent\begin{minipage}[t]{.37\linewidth}\vspace{0pt}\centering
\begin{tabular}{cc}\toprule
$\cTwoValue{x}$ (\cTwoValue{USD}) & $\cTwoMassHL{p_X(x)}$ \\ \midrule
$\cTwoValue{0}$ & $\cTwoMassHL{0.12}$ \\
$\cTwoValue{2}$ & $\cTwoMassHL{0.33}$ \\
$\cTwoValue{6}$ & $\cTwoMassHL{0.3}$ \\
$\cTwoValue{11}$ & $\cTwoMassHL{0.25}$ \\
\bottomrule\end{tabular}
\end{minipage}\hfill\begin{minipage}[t]{.60\linewidth}\vspace{0pt}\centering
\GTenPMFGraph[color=cTwoMassColor!20,xhighlight=GTenHLValue,yhighlight=cTwoMassMarker,width=\linewidth,height=4.15cm,xlabel={$\cTwoValue{x}$ (\cTwoValue{USD})},ylabel={$\cTwoMassHL{p_X(x)}$},xstep=1]{0,2,6,11}{0.12,0.33,0.3,0.25}
\end{minipage}\par
\textbf{PMF check.} All masses are positive and $0.12+0.33+0.3+0.25=1$. The graph uses the actual numerical spacing; gaps carry no mass.
\task{2}{Accumulate in increasing support order: $0.12;\ 0.45;\ 0.75;\ 1$. Each new height is the previous height plus the next PMF mass.}
\par\noindent\begin{minipage}[t]{.37\linewidth}\vspace{0pt}\centering
\begin{tabular}{cc}\toprule
Cutoff interval & $\cTwoCDFHL{F_X(x)}$ \\ \midrule
$x<0$ & $\cTwoCDFHL{0}$ \\
$0\leq x<2$ & $\cTwoCDFHL{0.12}$ \\
$2\leq x<6$ & $\cTwoCDFHL{0.45}$ \\
$6\leq x<11$ & $\cTwoCDFHL{0.75}$ \\
$x\geq 11$ & $\cTwoCDFHL{1}$ \\
\bottomrule\end{tabular}
\end{minipage}\hfill\begin{minipage}[t]{.60\linewidth}\vspace{0pt}\centering
\begin{tikzpicture}\begin{axis}[classCDF,xmin=-2,xmax=13,xtick={0,2,6,11}]
\addplot[cTwoCDFColor,thick,<-] coordinates {(-2,0)(0,0)};
\addplot[cTwoCDFColor,thick] coordinates {(0,0.12)(2,0.12)};
\addplot[cTwoCDFColor,thick] coordinates {(2,0.45)(6,0.45)};
\addplot[cTwoCDFColor,thick] coordinates {(6,0.75)(11,0.75)};
\addplot[cTwoCDFColor,thick,->] coordinates {(11,1)(13,1)};
\addplot[only marks,mark=*,cTwoCDFColor,mark size=1.8pt,mark options={fill=white}] coordinates {(0,0)(2,0.12)(6,0.45)(11,0.75)};
\addplot[only marks,mark=*,cTwoCDFColor,mark size=1.8pt] coordinates {(0,0.12)(2,0.45)(6,0.75)(11,1)};
\end{axis}\end{tikzpicture}
\end{minipage}\par
\textbf{CDF check.} All heights lie in $[0,1]$ and never decrease. Filled left/open right endpoints make the graph right-continuous. Horizontal tails give limits 0 and 1; the CDF reaches 1 at 11 USD. The jumps, $0.12,\ 0.33,\ 0.3,\ 0.25$, reproduce the PMF.
\textbf{Interpretation.} $F_X(6)=0.75$: the probability that the buyback payment is at most 6 USD (every type except aluminium) is 0.75.
\vfill
\clearpage
\problemtitle{10}{Device repair claims}{D}
\textbf{Missing probability and capped payout.} A device repair fund can receive only 0, 1, 2, 3 or 4 claims in one day. Their probabilities, in that order, are .30, .15, a missing probability, .20, .15. It pays 4 USD per claim, capped at 12 USD in total for the day; it pays the smaller of 4 times the count and 12 USD. There are no other payments. Record the daily payout.
\task{1}{Let $\cTwoValue{X}$ be the total payout for one day, in USD. The missing probability is $1-(.30+.15+.20+.15)=.20$. Apply $X=\min(4N,12)$ to the claim count $N$: payouts are $0,4,8,12,12$. Four claims would pay 16 USD without the cap. Group the last two categories.}
\begin{center}
\begin{tabular}{ccc}\toprule
Complete outcome & Probability & $\cTwoValue{X}$ (\cTwoValue{USD}) \\ \midrule
\cTwoOutcome{0 claims} & $\cTwoMassHL{0.3}$ & $\cTwoValue{0}$ \\
\cTwoOutcome{1 claims} & $\cTwoMassHL{0.15}$ & $\cTwoValue{4}$ \\
\cTwoOutcome{2 claims} & $\cTwoMassHL{0.2}$ & $\cTwoValue{8}$ \\
\cTwoOutcome{3 claims} & $\cTwoMassHL{0.2}$ & $\cTwoValue{12}$ \\
\cTwoOutcome{4 claims} & $\cTwoMassHL{0.15}$ & $\cTwoValue{12}$ \\
\bottomrule\end{tabular}
\end{center}
$\cTwoValue{S_X}=\cTwoValue{\{0,4,8,12\}}$; $p_X(x)=0$ elsewhere. The grouped PMF is:
\par\noindent\begin{minipage}[t]{.37\linewidth}\vspace{0pt}\centering
\begin{tabular}{cc}\toprule
$\cTwoValue{x}$ (\cTwoValue{USD}) & $\cTwoMassHL{p_X(x)}$ \\ \midrule
$\cTwoValue{0}$ & $\cTwoMassHL{0.3}$ \\
$\cTwoValue{4}$ & $\cTwoMassHL{0.15}$ \\
$\cTwoValue{8}$ & $\cTwoMassHL{0.2}$ \\
$\cTwoValue{12}$ & $\cTwoMassHL{0.35}$ \\
\bottomrule\end{tabular}
\end{minipage}\hfill\begin{minipage}[t]{.60\linewidth}\vspace{0pt}\centering
\GTenPMFGraph[color=cTwoMassColor!20,xhighlight=GTenHLValue,yhighlight=cTwoMassMarker,width=\linewidth,height=4.15cm,xlabel={$\cTwoValue{x}$ (\cTwoValue{USD})},ylabel={$\cTwoMassHL{p_X(x)}$},xstep=1]{0,4,8,12}{0.3,0.15,0.2,0.35}
\end{minipage}\par
\textbf{PMF check.} All masses are positive and $0.3+0.15+0.2+0.35=1$. The graph uses the actual numerical spacing; gaps carry no mass.
\task{2}{Accumulate in increasing support order: $0.3;\ 0.45;\ 0.65;\ 1$. Each new height is the previous height plus the next PMF mass.}
\par\noindent\begin{minipage}[t]{.37\linewidth}\vspace{0pt}\centering
\begin{tabular}{cc}\toprule
Cutoff interval & $\cTwoCDFHL{F_X(x)}$ \\ \midrule
$x<0$ & $\cTwoCDFHL{0}$ \\
$0\leq x<4$ & $\cTwoCDFHL{0.3}$ \\
$4\leq x<8$ & $\cTwoCDFHL{0.45}$ \\
$8\leq x<12$ & $\cTwoCDFHL{0.65}$ \\
$x\geq 12$ & $\cTwoCDFHL{1}$ \\
\bottomrule\end{tabular}
\end{minipage}\hfill\begin{minipage}[t]{.60\linewidth}\vspace{0pt}\centering
\begin{tikzpicture}\begin{axis}[classCDF,xmin=-4,xmax=16,xtick={0,4,8,12}]
\addplot[cTwoCDFColor,thick,<-] coordinates {(-4,0)(0,0)};
\addplot[cTwoCDFColor,thick] coordinates {(0,0.3)(4,0.3)};
\addplot[cTwoCDFColor,thick] coordinates {(4,0.45)(8,0.45)};
\addplot[cTwoCDFColor,thick] coordinates {(8,0.65)(12,0.65)};
\addplot[cTwoCDFColor,thick,->] coordinates {(12,1)(16,1)};
\addplot[only marks,mark=*,cTwoCDFColor,mark size=1.8pt,mark options={fill=white}] coordinates {(0,0)(4,0.3)(8,0.45)(12,0.65)};
\addplot[only marks,mark=*,cTwoCDFColor,mark size=1.8pt] coordinates {(0,0.3)(4,0.45)(8,0.65)(12,1)};
\end{axis}\end{tikzpicture}
\end{minipage}\par
\textbf{CDF check.} All heights lie in $[0,1]$ and never decrease. Filled left/open right endpoints make the graph right-continuous. Horizontal tails give limits 0 and 1; the CDF reaches 1 at 12 USD. The jumps, $0.3,\ 0.15,\ 0.2,\ 0.35$, reproduce the PMF.
\textbf{Interpretation.} $F_X(8)=0.65$: the probability that the payout is at most 8 USD (at most two claims) is 0.65.
\vfill
\clearpage
\problemtitle{11}{Mug sales}{E}
\textbf{Independent purchases and discount.} A school enterprise offers one mug to each of three visitors. Independently, each buys one with probability 0.6 or passes with probability 0.4, regardless of the discount. Each item costs 6 USD. Only if all three buy, subtract a 4 USD discount from the combined bill. Otherwise there is no discount. There are no other sales or charges. Record total revenue.
\task{1}{Let $\cTwoValue{X}$ be the total revenue from the three visitors, in USD. Use B for buying and N for passing, in visitor order. Multiply within each complete sequence: e.g. $P(\mathrm{BNN})=0.6(0.4)^2=0.096$. The three one-buyer and three two-buyer outcomes each form a group. Only BBB receives the discount: $3(6)-4=14$ USD.}
\begin{center}
\begin{tabular}{ccc}\toprule
Complete outcome & Probability & $\cTwoValue{X}$ (\cTwoValue{USD}) \\ \midrule
\cTwoOutcome{NNN} & $\cTwoMassHL{0.064}$ & $\cTwoValue{0}$ \\
\cTwoOutcome{BNN} & $\cTwoMassHL{0.096}$ & $\cTwoValue{6}$ \\
\cTwoOutcome{NBN} & $\cTwoMassHL{0.096}$ & $\cTwoValue{6}$ \\
\cTwoOutcome{NNB} & $\cTwoMassHL{0.096}$ & $\cTwoValue{6}$ \\
\cTwoOutcome{BBN} & $\cTwoMassHL{0.144}$ & $\cTwoValue{12}$ \\
\cTwoOutcome{BNB} & $\cTwoMassHL{0.144}$ & $\cTwoValue{12}$ \\
\cTwoOutcome{NBB} & $\cTwoMassHL{0.144}$ & $\cTwoValue{12}$ \\
\cTwoOutcome{BBB} & $\cTwoMassHL{0.216}$ & $\cTwoValue{14}$ \\
\bottomrule\end{tabular}
\end{center}
$\cTwoValue{S_X}=\cTwoValue{\{0,6,12,14\}}$; $p_X(x)=0$ elsewhere. The grouped PMF is:
\par\noindent\begin{minipage}[t]{.37\linewidth}\vspace{0pt}\centering
\begin{tabular}{cc}\toprule
$\cTwoValue{x}$ (\cTwoValue{USD}) & $\cTwoMassHL{p_X(x)}$ \\ \midrule
$\cTwoValue{0}$ & $\cTwoMassHL{0.064}$ \\
$\cTwoValue{6}$ & $\cTwoMassHL{0.288}$ \\
$\cTwoValue{12}$ & $\cTwoMassHL{0.432}$ \\
$\cTwoValue{14}$ & $\cTwoMassHL{0.216}$ \\
\bottomrule\end{tabular}
\end{minipage}\hfill\begin{minipage}[t]{.60\linewidth}\vspace{0pt}\centering
\GTenPMFGraph[color=cTwoMassColor!20,xhighlight=GTenHLValue,yhighlight=cTwoMassMarker,width=\linewidth,height=4.15cm,xlabel={$\cTwoValue{x}$ (\cTwoValue{USD})},ylabel={$\cTwoMassHL{p_X(x)}$},xstep=1]{0,6,12,14}{0.064,0.288,0.432,0.216}
\end{minipage}\par
\textbf{PMF check.} All masses are positive and $0.064+0.288+0.432+0.216=1$. The graph uses the actual numerical spacing; gaps carry no mass.
\task{2}{Accumulate in increasing support order: $0.064;\ 0.352;\ 0.784;\ 1$. Each new height is the previous height plus the next PMF mass.}
\par\noindent\begin{minipage}[t]{.37\linewidth}\vspace{0pt}\centering
\begin{tabular}{cc}\toprule
Cutoff interval & $\cTwoCDFHL{F_X(x)}$ \\ \midrule
$x<0$ & $\cTwoCDFHL{0}$ \\
$0\leq x<6$ & $\cTwoCDFHL{0.064}$ \\
$6\leq x<12$ & $\cTwoCDFHL{0.352}$ \\
$12\leq x<14$ & $\cTwoCDFHL{0.784}$ \\
$x\geq 14$ & $\cTwoCDFHL{1}$ \\
\bottomrule\end{tabular}
\end{minipage}\hfill\begin{minipage}[t]{.60\linewidth}\vspace{0pt}\centering
\begin{tikzpicture}\begin{axis}[classCDF,xmin=-2,xmax=16,xtick={0,6,12,14}]
\addplot[cTwoCDFColor,thick,<-] coordinates {(-2,0)(0,0)};
\addplot[cTwoCDFColor,thick] coordinates {(0,0.064)(6,0.064)};
\addplot[cTwoCDFColor,thick] coordinates {(6,0.352)(12,0.352)};
\addplot[cTwoCDFColor,thick] coordinates {(12,0.784)(14,0.784)};
\addplot[cTwoCDFColor,thick,->] coordinates {(14,1)(16,1)};
\addplot[only marks,mark=*,cTwoCDFColor,mark size=1.8pt,mark options={fill=white}] coordinates {(0,0)(6,0.064)(12,0.352)(14,0.784)};
\addplot[only marks,mark=*,cTwoCDFColor,mark size=1.8pt] coordinates {(0,0.064)(6,0.352)(12,0.784)(14,1)};
\end{axis}\end{tikzpicture}
\end{minipage}\par
\textbf{CDF check.} All heights lie in $[0,1]$ and never decrease. Filled left/open right endpoints make the graph right-continuous. Horizontal tails give limits 0 and 1; the CDF reaches 1 at 14 USD. The jumps, $0.064,\ 0.288,\ 0.432,\ 0.216$, reproduce the PMF.
\textbf{Interpretation.} $F_X(12)=0.784$: the probability that revenue is at most 12 USD (not all three visitors buy) is 0.784.
\vfill
\clearpage
\problemtitle{12}{Two-stage garden fund}{F}
\textbf{Signed two-stage net results.} A hypothetical classroom fund receives two cash returns. The first is 12 USD, 0 USD or 5 USD with probabilities .40, .25, .35, respectively. Independently, the second is 0 USD with probability 0.7 or 5 USD with probability 0.3. These are all possible returns. A fixed participation cost of 5 USD is charged once. There are no other receipts or charges. Record the sum of both returns minus the cost; a negative net result means a loss.
\task{1}{Let $\cTwoValue{X}$ be the signed net result after both stages and the cost, in USD. For each ordered pair, multiply independent probabilities, then compute $X=R_1+R_2-5$. For example, $(12,0)$ has probability $.40(0.7)=0.28$ and net result $12+0-5=7$. The pairs $(0,5)$ and $(5,0)$ both break even; group their masses. Subtract the cost only once and sort signed results numerically.}
\begin{center}
\begin{tabular}{ccc}\toprule
Complete outcome & Probability & $\cTwoValue{X}$ (\cTwoValue{USD}) \\ \midrule
\cTwoOutcome{$(12,0)$} & $\cTwoMassHL{0.28}$ & $\cTwoValue{7}$ \\
\cTwoOutcome{$(12,5)$} & $\cTwoMassHL{0.12}$ & $\cTwoValue{12}$ \\
\cTwoOutcome{$(0,0)$} & $\cTwoMassHL{0.175}$ & $\cTwoValue{-5}$ \\
\cTwoOutcome{$(0,5)$} & $\cTwoMassHL{0.075}$ & $\cTwoValue{0}$ \\
\cTwoOutcome{$(5,0)$} & $\cTwoMassHL{0.245}$ & $\cTwoValue{0}$ \\
\cTwoOutcome{$(5,5)$} & $\cTwoMassHL{0.105}$ & $\cTwoValue{5}$ \\
\bottomrule\end{tabular}
\end{center}
$\cTwoValue{S_X}=\cTwoValue{\{-5,0,5,7,12\}}$; $p_X(x)=0$ elsewhere. The grouped PMF is:
\par\noindent\begin{minipage}[t]{.37\linewidth}\vspace{0pt}\centering
\begin{tabular}{cc}\toprule
$\cTwoValue{x}$ (\cTwoValue{USD}) & $\cTwoMassHL{p_X(x)}$ \\ \midrule
$\cTwoValue{-5}$ & $\cTwoMassHL{0.175}$ \\
$\cTwoValue{0}$ & $\cTwoMassHL{0.32}$ \\
$\cTwoValue{5}$ & $\cTwoMassHL{0.105}$ \\
$\cTwoValue{7}$ & $\cTwoMassHL{0.28}$ \\
$\cTwoValue{12}$ & $\cTwoMassHL{0.12}$ \\
\bottomrule\end{tabular}
\end{minipage}\hfill\begin{minipage}[t]{.60\linewidth}\vspace{0pt}\centering
\GTenPMFGraph[color=cTwoMassColor!20,xhighlight=GTenHLValue,yhighlight=cTwoMassMarker,width=\linewidth,height=4.15cm,xlabel={$\cTwoValue{x}$ (\cTwoValue{USD})},ylabel={$\cTwoMassHL{p_X(x)}$},xstep=1]{-5,0,5,7,12}{0.175,0.32,0.105,0.28,0.12}
\end{minipage}\par
\textbf{PMF check.} All masses are positive and $0.175+0.32+0.105+0.28+0.12=1$. The graph uses the actual numerical spacing; gaps carry no mass.
\task{2}{Accumulate in increasing support order: $0.175;\ 0.495;\ 0.6;\ 0.88;\ 1$. Each new height is the previous height plus the next PMF mass.}
\par\noindent\begin{minipage}[t]{.37\linewidth}\vspace{0pt}\centering
\begin{tabular}{cc}\toprule
Cutoff interval & $\cTwoCDFHL{F_X(x)}$ \\ \midrule
$x<-5$ & $\cTwoCDFHL{0}$ \\
$-5\leq x<0$ & $\cTwoCDFHL{0.175}$ \\
$0\leq x<5$ & $\cTwoCDFHL{0.495}$ \\
$5\leq x<7$ & $\cTwoCDFHL{0.6}$ \\
$7\leq x<12$ & $\cTwoCDFHL{0.88}$ \\
$x\geq 12$ & $\cTwoCDFHL{1}$ \\
\bottomrule\end{tabular}
\end{minipage}\hfill\begin{minipage}[t]{.60\linewidth}\vspace{0pt}\centering
\begin{tikzpicture}\begin{axis}[classCDF,xmin=-7,xmax=14,xtick={-5,0,5,7,12}]
\addplot[cTwoCDFColor,thick,<-] coordinates {(-7,0)(-5,0)};
\addplot[cTwoCDFColor,thick] coordinates {(-5,0.175)(0,0.175)};
\addplot[cTwoCDFColor,thick] coordinates {(0,0.495)(5,0.495)};
\addplot[cTwoCDFColor,thick] coordinates {(5,0.6)(7,0.6)};
\addplot[cTwoCDFColor,thick] coordinates {(7,0.88)(12,0.88)};
\addplot[cTwoCDFColor,thick,->] coordinates {(12,1)(14,1)};
\addplot[only marks,mark=*,cTwoCDFColor,mark size=1.8pt,mark options={fill=white}] coordinates {(-5,0)(0,0.175)(5,0.495)(7,0.6)(12,0.88)};
\addplot[only marks,mark=*,cTwoCDFColor,mark size=1.8pt] coordinates {(-5,0.175)(0,0.495)(5,0.6)(7,0.88)(12,1)};
\end{axis}\end{tikzpicture}
\end{minipage}\par
\textbf{CDF check.} All heights lie in $[0,1]$ and never decrease. Filled left/open right endpoints make the graph right-continuous. Horizontal tails give limits 0 and 1; the CDF reaches 1 at 12 USD. The jumps, $0.175,\ 0.32,\ 0.105,\ 0.28,\ 0.12$, reproduce the PMF.
\textbf{Interpretation.} $F_X(0)=0.495$: the probability that the net result is nonpositive (a loss or break-even) is 0.495.
\vfill
\clearpage
\problemtitle{13}{Theatre ticket payments}{A}
\textbf{Payment fees.} A theatre ticket booth receives one payment. The only methods are cash, debit card, credit card and mobile wallet, with probabilities .30, .20, .15, .35, respectively. Cash has no fee; either card costs 4 USD; a mobile wallet costs 8 USD. There are no other charges. Record the processing fee.
\task{1}{Let $\cTwoValue{X}$ be the processing fee for one payment, in USD. Cash maps to 0, either card to 4, and wallet to 8. The two card events are mutually exclusive, so their masses add: $p_X(4)=.20+.15=0.35$.}
\begin{center}
\begin{tabular}{ccc}\toprule
Complete outcome & Probability & $\cTwoValue{X}$ (\cTwoValue{USD}) \\ \midrule
\cTwoOutcome{Cash} & $\cTwoMassHL{0.3}$ & $\cTwoValue{0}$ \\
\cTwoOutcome{Debit card} & $\cTwoMassHL{0.2}$ & $\cTwoValue{4}$ \\
\cTwoOutcome{Credit card} & $\cTwoMassHL{0.15}$ & $\cTwoValue{4}$ \\
\cTwoOutcome{Wallet} & $\cTwoMassHL{0.35}$ & $\cTwoValue{8}$ \\
\bottomrule\end{tabular}
\end{center}
$\cTwoValue{S_X}=\cTwoValue{\{0,4,8\}}$; $p_X(x)=0$ elsewhere. The grouped PMF is:
\par\noindent\begin{minipage}[t]{.37\linewidth}\vspace{0pt}\centering
\begin{tabular}{cc}\toprule
$\cTwoValue{x}$ (\cTwoValue{USD}) & $\cTwoMassHL{p_X(x)}$ \\ \midrule
$\cTwoValue{0}$ & $\cTwoMassHL{0.3}$ \\
$\cTwoValue{4}$ & $\cTwoMassHL{0.35}$ \\
$\cTwoValue{8}$ & $\cTwoMassHL{0.35}$ \\
\bottomrule\end{tabular}
\end{minipage}\hfill\begin{minipage}[t]{.60\linewidth}\vspace{0pt}\centering
\GTenPMFGraph[color=cTwoMassColor!20,xhighlight=GTenHLValue,yhighlight=cTwoMassMarker,width=\linewidth,height=4.15cm,xlabel={$\cTwoValue{x}$ (\cTwoValue{USD})},ylabel={$\cTwoMassHL{p_X(x)}$},xstep=1]{0,4,8}{0.3,0.35,0.35}
\end{minipage}\par
\textbf{PMF check.} All masses are positive and $0.3+0.35+0.35=1$. The graph uses the actual numerical spacing; gaps carry no mass.
\task{2}{Accumulate in increasing support order: $0.3;\ 0.65;\ 1$. Each new height is the previous height plus the next PMF mass.}
\par\noindent\begin{minipage}[t]{.37\linewidth}\vspace{0pt}\centering
\begin{tabular}{cc}\toprule
Cutoff interval & $\cTwoCDFHL{F_X(x)}$ \\ \midrule
$x<0$ & $\cTwoCDFHL{0}$ \\
$0\leq x<4$ & $\cTwoCDFHL{0.3}$ \\
$4\leq x<8$ & $\cTwoCDFHL{0.65}$ \\
$x\geq 8$ & $\cTwoCDFHL{1}$ \\
\bottomrule\end{tabular}
\end{minipage}\hfill\begin{minipage}[t]{.60\linewidth}\vspace{0pt}\centering
\begin{tikzpicture}\begin{axis}[classCDF,xmin=-4,xmax=12,xtick={0,4,8}]
\addplot[cTwoCDFColor,thick,<-] coordinates {(-4,0)(0,0)};
\addplot[cTwoCDFColor,thick] coordinates {(0,0.3)(4,0.3)};
\addplot[cTwoCDFColor,thick] coordinates {(4,0.65)(8,0.65)};
\addplot[cTwoCDFColor,thick,->] coordinates {(8,1)(12,1)};
\addplot[only marks,mark=*,cTwoCDFColor,mark size=1.8pt,mark options={fill=white}] coordinates {(0,0)(4,0.3)(8,0.65)};
\addplot[only marks,mark=*,cTwoCDFColor,mark size=1.8pt] coordinates {(0,0.3)(4,0.65)(8,1)};
\end{axis}\end{tikzpicture}
\end{minipage}\par
\textbf{CDF check.} All heights lie in $[0,1]$ and never decrease. Filled left/open right endpoints make the graph right-continuous. Horizontal tails give limits 0 and 1; the CDF reaches 1 at 8 USD. The jumps, $0.3,\ 0.35,\ 0.35$, reproduce the PMF.
\textbf{Interpretation.} $F_X(4)=0.65$: the probability that the fee is at most 4 USD (cash or either card) is 0.65.
\vfill
\clearpage
\problemtitle{14}{Sports stall payments}{A}
\textbf{Payment fees.} A sports stall receives one payment. The only methods are cash, debit card, credit card and mobile wallet, with probabilities .12, .28, .22, .38, respectively. Cash has no fee; either card costs 5 USD; a mobile wallet costs 10 USD. There are no other charges. Record the processing fee.
\task{1}{Let $\cTwoValue{X}$ be the processing fee for one payment, in USD. Cash maps to 0, either card to 5, and wallet to 10. The two card events are mutually exclusive, so their masses add: $p_X(5)=.28+.22=0.5$.}
\begin{center}
\begin{tabular}{ccc}\toprule
Complete outcome & Probability & $\cTwoValue{X}$ (\cTwoValue{USD}) \\ \midrule
\cTwoOutcome{Cash} & $\cTwoMassHL{0.12}$ & $\cTwoValue{0}$ \\
\cTwoOutcome{Debit card} & $\cTwoMassHL{0.28}$ & $\cTwoValue{5}$ \\
\cTwoOutcome{Credit card} & $\cTwoMassHL{0.22}$ & $\cTwoValue{5}$ \\
\cTwoOutcome{Wallet} & $\cTwoMassHL{0.38}$ & $\cTwoValue{10}$ \\
\bottomrule\end{tabular}
\end{center}
$\cTwoValue{S_X}=\cTwoValue{\{0,5,10\}}$; $p_X(x)=0$ elsewhere. The grouped PMF is:
\par\noindent\begin{minipage}[t]{.37\linewidth}\vspace{0pt}\centering
\begin{tabular}{cc}\toprule
$\cTwoValue{x}$ (\cTwoValue{USD}) & $\cTwoMassHL{p_X(x)}$ \\ \midrule
$\cTwoValue{0}$ & $\cTwoMassHL{0.12}$ \\
$\cTwoValue{5}$ & $\cTwoMassHL{0.5}$ \\
$\cTwoValue{10}$ & $\cTwoMassHL{0.38}$ \\
\bottomrule\end{tabular}
\end{minipage}\hfill\begin{minipage}[t]{.60\linewidth}\vspace{0pt}\centering
\GTenPMFGraph[color=cTwoMassColor!20,xhighlight=GTenHLValue,yhighlight=cTwoMassMarker,width=\linewidth,height=4.15cm,xlabel={$\cTwoValue{x}$ (\cTwoValue{USD})},ylabel={$\cTwoMassHL{p_X(x)}$},xstep=1]{0,5,10}{0.12,0.5,0.38}
\end{minipage}\par
\textbf{PMF check.} All masses are positive and $0.12+0.5+0.38=1$. The graph uses the actual numerical spacing; gaps carry no mass.
\task{2}{Accumulate in increasing support order: $0.12;\ 0.62;\ 1$. Each new height is the previous height plus the next PMF mass.}
\par\noindent\begin{minipage}[t]{.37\linewidth}\vspace{0pt}\centering
\begin{tabular}{cc}\toprule
Cutoff interval & $\cTwoCDFHL{F_X(x)}$ \\ \midrule
$x<0$ & $\cTwoCDFHL{0}$ \\
$0\leq x<5$ & $\cTwoCDFHL{0.12}$ \\
$5\leq x<10$ & $\cTwoCDFHL{0.62}$ \\
$x\geq 10$ & $\cTwoCDFHL{1}$ \\
\bottomrule\end{tabular}
\end{minipage}\hfill\begin{minipage}[t]{.60\linewidth}\vspace{0pt}\centering
\begin{tikzpicture}\begin{axis}[classCDF,xmin=-5,xmax=15,xtick={0,5,10}]
\addplot[cTwoCDFColor,thick,<-] coordinates {(-5,0)(0,0)};
\addplot[cTwoCDFColor,thick] coordinates {(0,0.12)(5,0.12)};
\addplot[cTwoCDFColor,thick] coordinates {(5,0.62)(10,0.62)};
\addplot[cTwoCDFColor,thick,->] coordinates {(10,1)(15,1)};
\addplot[only marks,mark=*,cTwoCDFColor,mark size=1.8pt,mark options={fill=white}] coordinates {(0,0)(5,0.12)(10,0.62)};
\addplot[only marks,mark=*,cTwoCDFColor,mark size=1.8pt] coordinates {(0,0.12)(5,0.62)(10,1)};
\end{axis}\end{tikzpicture}
\end{minipage}\par
\textbf{CDF check.} All heights lie in $[0,1]$ and never decrease. Filled left/open right endpoints make the graph right-continuous. Horizontal tails give limits 0 and 1; the CDF reaches 1 at 10 USD. The jumps, $0.12,\ 0.5,\ 0.38$, reproduce the PMF.
\textbf{Interpretation.} $F_X(5)=0.62$: the probability that the fee is at most 5 USD (cash or either card) is 0.62.
\vfill
\clearpage
\problemtitle{15}{Camera returns}{B}
\textbf{Late fees.} A rental service records the return of one camera. The exhaustive categories are on time or early, 1--5 whole minutes late, 6--15 whole minutes late, and more than 15 minutes late, with probabilities .18, .22, .35, .25, respectively. The late fee is 0 USD, 3 USD, 3 USD and 6 USD, respectively. No other charges apply. Record the late fee.
\task{1}{Let $\cTwoValue{X}$ be the late fee for one return, in USD. The two shorter delays both map to 3 USD. Their mutually exclusive masses add: $p_X(3)=.22+.35=0.57$. On-time returns map to 0 and the longest delay maps to 6.}
\begin{center}
\begin{tabular}{ccc}\toprule
Complete outcome & Probability & $\cTwoValue{X}$ (\cTwoValue{USD}) \\ \midrule
\cTwoOutcome{On time/early} & $\cTwoMassHL{0.18}$ & $\cTwoValue{0}$ \\
\cTwoOutcome{1--5 min} & $\cTwoMassHL{0.22}$ & $\cTwoValue{3}$ \\
\cTwoOutcome{6--15 min} & $\cTwoMassHL{0.35}$ & $\cTwoValue{3}$ \\
\cTwoOutcome{Over 15 min} & $\cTwoMassHL{0.25}$ & $\cTwoValue{6}$ \\
\bottomrule\end{tabular}
\end{center}
$\cTwoValue{S_X}=\cTwoValue{\{0,3,6\}}$; $p_X(x)=0$ elsewhere. The grouped PMF is:
\par\noindent\begin{minipage}[t]{.37\linewidth}\vspace{0pt}\centering
\begin{tabular}{cc}\toprule
$\cTwoValue{x}$ (\cTwoValue{USD}) & $\cTwoMassHL{p_X(x)}$ \\ \midrule
$\cTwoValue{0}$ & $\cTwoMassHL{0.18}$ \\
$\cTwoValue{3}$ & $\cTwoMassHL{0.57}$ \\
$\cTwoValue{6}$ & $\cTwoMassHL{0.25}$ \\
\bottomrule\end{tabular}
\end{minipage}\hfill\begin{minipage}[t]{.60\linewidth}\vspace{0pt}\centering
\GTenPMFGraph[color=cTwoMassColor!20,xhighlight=GTenHLValue,yhighlight=cTwoMassMarker,width=\linewidth,height=4.15cm,xlabel={$\cTwoValue{x}$ (\cTwoValue{USD})},ylabel={$\cTwoMassHL{p_X(x)}$},xstep=1]{0,3,6}{0.18,0.57,0.25}
\end{minipage}\par
\textbf{PMF check.} All masses are positive and $0.18+0.57+0.25=1$. The graph uses the actual numerical spacing; gaps carry no mass.
\task{2}{Accumulate in increasing support order: $0.18;\ 0.75;\ 1$. Each new height is the previous height plus the next PMF mass.}
\par\noindent\begin{minipage}[t]{.37\linewidth}\vspace{0pt}\centering
\begin{tabular}{cc}\toprule
Cutoff interval & $\cTwoCDFHL{F_X(x)}$ \\ \midrule
$x<0$ & $\cTwoCDFHL{0}$ \\
$0\leq x<3$ & $\cTwoCDFHL{0.18}$ \\
$3\leq x<6$ & $\cTwoCDFHL{0.75}$ \\
$x\geq 6$ & $\cTwoCDFHL{1}$ \\
\bottomrule\end{tabular}
\end{minipage}\hfill\begin{minipage}[t]{.60\linewidth}\vspace{0pt}\centering
\begin{tikzpicture}\begin{axis}[classCDF,xmin=-3,xmax=9,xtick={0,3,6}]
\addplot[cTwoCDFColor,thick,<-] coordinates {(-3,0)(0,0)};
\addplot[cTwoCDFColor,thick] coordinates {(0,0.18)(3,0.18)};
\addplot[cTwoCDFColor,thick] coordinates {(3,0.75)(6,0.75)};
\addplot[cTwoCDFColor,thick,->] coordinates {(6,1)(9,1)};
\addplot[only marks,mark=*,cTwoCDFColor,mark size=1.8pt,mark options={fill=white}] coordinates {(0,0)(3,0.18)(6,0.75)};
\addplot[only marks,mark=*,cTwoCDFColor,mark size=1.8pt] coordinates {(0,0.18)(3,0.75)(6,1)};
\end{axis}\end{tikzpicture}
\end{minipage}\par
\textbf{CDF check.} All heights lie in $[0,1]$ and never decrease. Filled left/open right endpoints make the graph right-continuous. Horizontal tails give limits 0 and 1; the CDF reaches 1 at 6 USD. The jumps, $0.18,\ 0.57,\ 0.25$, reproduce the PMF.
\textbf{Interpretation.} $F_X(3)=0.75$: the probability that the late fee is at most 3 USD (the return is no more than 15 minutes late) is 0.75.
\vfill
\clearpage
\problemtitle{16}{Tool returns}{B}
\textbf{Late fees.} A rental service records the return of one tool. The exhaustive categories are on time or early, 1--5 whole minutes late, 6--15 whole minutes late, and more than 15 minutes late, with probabilities .35, .15, .20, .30, respectively. The late fee is 0 USD, 5 USD, 5 USD and 10 USD, respectively. No other charges apply. Record the late fee.
\task{1}{Let $\cTwoValue{X}$ be the late fee for one return, in USD. The two shorter delays both map to 5 USD. Their mutually exclusive masses add: $p_X(5)=.15+.20=0.35$. On-time returns map to 0 and the longest delay maps to 10.}
\begin{center}
\begin{tabular}{ccc}\toprule
Complete outcome & Probability & $\cTwoValue{X}$ (\cTwoValue{USD}) \\ \midrule
\cTwoOutcome{On time/early} & $\cTwoMassHL{0.35}$ & $\cTwoValue{0}$ \\
\cTwoOutcome{1--5 min} & $\cTwoMassHL{0.15}$ & $\cTwoValue{5}$ \\
\cTwoOutcome{6--15 min} & $\cTwoMassHL{0.2}$ & $\cTwoValue{5}$ \\
\cTwoOutcome{Over 15 min} & $\cTwoMassHL{0.3}$ & $\cTwoValue{10}$ \\
\bottomrule\end{tabular}
\end{center}
$\cTwoValue{S_X}=\cTwoValue{\{0,5,10\}}$; $p_X(x)=0$ elsewhere. The grouped PMF is:
\par\noindent\begin{minipage}[t]{.37\linewidth}\vspace{0pt}\centering
\begin{tabular}{cc}\toprule
$\cTwoValue{x}$ (\cTwoValue{USD}) & $\cTwoMassHL{p_X(x)}$ \\ \midrule
$\cTwoValue{0}$ & $\cTwoMassHL{0.35}$ \\
$\cTwoValue{5}$ & $\cTwoMassHL{0.35}$ \\
$\cTwoValue{10}$ & $\cTwoMassHL{0.3}$ \\
\bottomrule\end{tabular}
\end{minipage}\hfill\begin{minipage}[t]{.60\linewidth}\vspace{0pt}\centering
\GTenPMFGraph[color=cTwoMassColor!20,xhighlight=GTenHLValue,yhighlight=cTwoMassMarker,width=\linewidth,height=4.15cm,xlabel={$\cTwoValue{x}$ (\cTwoValue{USD})},ylabel={$\cTwoMassHL{p_X(x)}$},xstep=1]{0,5,10}{0.35,0.35,0.3}
\end{minipage}\par
\textbf{PMF check.} All masses are positive and $0.35+0.35+0.3=1$. The graph uses the actual numerical spacing; gaps carry no mass.
\task{2}{Accumulate in increasing support order: $0.35;\ 0.7;\ 1$. Each new height is the previous height plus the next PMF mass.}
\par\noindent\begin{minipage}[t]{.37\linewidth}\vspace{0pt}\centering
\begin{tabular}{cc}\toprule
Cutoff interval & $\cTwoCDFHL{F_X(x)}$ \\ \midrule
$x<0$ & $\cTwoCDFHL{0}$ \\
$0\leq x<5$ & $\cTwoCDFHL{0.35}$ \\
$5\leq x<10$ & $\cTwoCDFHL{0.7}$ \\
$x\geq 10$ & $\cTwoCDFHL{1}$ \\
\bottomrule\end{tabular}
\end{minipage}\hfill\begin{minipage}[t]{.60\linewidth}\vspace{0pt}\centering
\begin{tikzpicture}\begin{axis}[classCDF,xmin=-5,xmax=15,xtick={0,5,10}]
\addplot[cTwoCDFColor,thick,<-] coordinates {(-5,0)(0,0)};
\addplot[cTwoCDFColor,thick] coordinates {(0,0.35)(5,0.35)};
\addplot[cTwoCDFColor,thick] coordinates {(5,0.7)(10,0.7)};
\addplot[cTwoCDFColor,thick,->] coordinates {(10,1)(15,1)};
\addplot[only marks,mark=*,cTwoCDFColor,mark size=1.8pt,mark options={fill=white}] coordinates {(0,0)(5,0.35)(10,0.7)};
\addplot[only marks,mark=*,cTwoCDFColor,mark size=1.8pt] coordinates {(0,0.35)(5,0.7)(10,1)};
\end{axis}\end{tikzpicture}
\end{minipage}\par
\textbf{CDF check.} All heights lie in $[0,1]$ and never decrease. Filled left/open right endpoints make the graph right-continuous. Horizontal tails give limits 0 and 1; the CDF reaches 1 at 10 USD. The jumps, $0.35,\ 0.35,\ 0.3$, reproduce the PMF.
\textbf{Interpretation.} $F_X(5)=0.7$: the probability that the late fee is at most 5 USD (the return is no more than 15 minutes late) is 0.7.
\vfill
\clearpage
\problemtitle{17}{Workshop scrap buyback}{C}
\textbf{Uneven buyback payments.} A workshop scrap centre buys one standard-sized bag containing only one material type. The exhaustive types are rejected material, paper, cardboard, plastic, steel and aluminium, with probabilities .20, .12, .18, .15, .10, .25, respectively. Payments are 0 USD for rejected material, 4 USD for either paper or cardboard, 9 USD for either plastic or steel, and 15 USD for aluminium. There are no other payments. Record the buyback payment.
\task{1}{Let $\cTwoValue{X}$ be the buyback payment for one bag, in USD. Paper and cardboard share 4 USD, so their probabilities add to 0.3. Plastic and steel share 9 USD, giving 0.25. Keep the actual gaps 4, 5 and 6 USD on both graphs.}
\begin{center}
\begin{tabular}{ccc}\toprule
Complete outcome & Probability & $\cTwoValue{X}$ (\cTwoValue{USD}) \\ \midrule
\cTwoOutcome{Rejected} & $\cTwoMassHL{0.2}$ & $\cTwoValue{0}$ \\
\cTwoOutcome{Paper} & $\cTwoMassHL{0.12}$ & $\cTwoValue{4}$ \\
\cTwoOutcome{Cardboard} & $\cTwoMassHL{0.18}$ & $\cTwoValue{4}$ \\
\cTwoOutcome{Plastic} & $\cTwoMassHL{0.15}$ & $\cTwoValue{9}$ \\
\cTwoOutcome{Steel} & $\cTwoMassHL{0.1}$ & $\cTwoValue{9}$ \\
\cTwoOutcome{Aluminium} & $\cTwoMassHL{0.25}$ & $\cTwoValue{15}$ \\
\bottomrule\end{tabular}
\end{center}
$\cTwoValue{S_X}=\cTwoValue{\{0,4,9,15\}}$; $p_X(x)=0$ elsewhere. The grouped PMF is:
\par\noindent\begin{minipage}[t]{.37\linewidth}\vspace{0pt}\centering
\begin{tabular}{cc}\toprule
$\cTwoValue{x}$ (\cTwoValue{USD}) & $\cTwoMassHL{p_X(x)}$ \\ \midrule
$\cTwoValue{0}$ & $\cTwoMassHL{0.2}$ \\
$\cTwoValue{4}$ & $\cTwoMassHL{0.3}$ \\
$\cTwoValue{9}$ & $\cTwoMassHL{0.25}$ \\
$\cTwoValue{15}$ & $\cTwoMassHL{0.25}$ \\
\bottomrule\end{tabular}
\end{minipage}\hfill\begin{minipage}[t]{.60\linewidth}\vspace{0pt}\centering
\GTenPMFGraph[color=cTwoMassColor!20,xhighlight=GTenHLValue,yhighlight=cTwoMassMarker,width=\linewidth,height=4.15cm,xlabel={$\cTwoValue{x}$ (\cTwoValue{USD})},ylabel={$\cTwoMassHL{p_X(x)}$},xstep=1]{0,4,9,15}{0.2,0.3,0.25,0.25}
\end{minipage}\par
\textbf{PMF check.} All masses are positive and $0.2+0.3+0.25+0.25=1$. The graph uses the actual numerical spacing; gaps carry no mass.
\task{2}{Accumulate in increasing support order: $0.2;\ 0.5;\ 0.75;\ 1$. Each new height is the previous height plus the next PMF mass.}
\par\noindent\begin{minipage}[t]{.37\linewidth}\vspace{0pt}\centering
\begin{tabular}{cc}\toprule
Cutoff interval & $\cTwoCDFHL{F_X(x)}$ \\ \midrule
$x<0$ & $\cTwoCDFHL{0}$ \\
$0\leq x<4$ & $\cTwoCDFHL{0.2}$ \\
$4\leq x<9$ & $\cTwoCDFHL{0.5}$ \\
$9\leq x<15$ & $\cTwoCDFHL{0.75}$ \\
$x\geq 15$ & $\cTwoCDFHL{1}$ \\
\bottomrule\end{tabular}
\end{minipage}\hfill\begin{minipage}[t]{.60\linewidth}\vspace{0pt}\centering
\begin{tikzpicture}\begin{axis}[classCDF,xmin=-4,xmax=19,xtick={0,4,9,15}]
\addplot[cTwoCDFColor,thick,<-] coordinates {(-4,0)(0,0)};
\addplot[cTwoCDFColor,thick] coordinates {(0,0.2)(4,0.2)};
\addplot[cTwoCDFColor,thick] coordinates {(4,0.5)(9,0.5)};
\addplot[cTwoCDFColor,thick] coordinates {(9,0.75)(15,0.75)};
\addplot[cTwoCDFColor,thick,->] coordinates {(15,1)(19,1)};
\addplot[only marks,mark=*,cTwoCDFColor,mark size=1.8pt,mark options={fill=white}] coordinates {(0,0)(4,0.2)(9,0.5)(15,0.75)};
\addplot[only marks,mark=*,cTwoCDFColor,mark size=1.8pt] coordinates {(0,0.2)(4,0.5)(9,0.75)(15,1)};
\end{axis}\end{tikzpicture}
\end{minipage}\par
\textbf{CDF check.} All heights lie in $[0,1]$ and never decrease. Filled left/open right endpoints make the graph right-continuous. Horizontal tails give limits 0 and 1; the CDF reaches 1 at 15 USD. The jumps, $0.2,\ 0.3,\ 0.25,\ 0.25$, reproduce the PMF.
\textbf{Interpretation.} $F_X(9)=0.75$: the probability that the buyback payment is at most 9 USD (every type except aluminium) is 0.75.
\vfill
\clearpage
\problemtitle{18}{Community salvage buyback}{C}
\textbf{Uneven buyback payments.} A community salvage stall buys one standard-sized bag containing only one material type. The exhaustive types are rejected material, paper, cardboard, plastic, steel and aluminium, with probabilities .08, .17, .15, .25, .15, .20, respectively. Payments are 0 USD for rejected material, 3 USD for either paper or cardboard, 8 USD for either plastic or steel, and 14 USD for aluminium. There are no other payments. Record the buyback payment.
\task{1}{Let $\cTwoValue{X}$ be the buyback payment for one bag, in USD. Paper and cardboard share 3 USD, so their probabilities add to 0.32. Plastic and steel share 8 USD, giving 0.4. Keep the actual gaps 3, 5 and 6 USD on both graphs.}
\begin{center}
\begin{tabular}{ccc}\toprule
Complete outcome & Probability & $\cTwoValue{X}$ (\cTwoValue{USD}) \\ \midrule
\cTwoOutcome{Rejected} & $\cTwoMassHL{0.08}$ & $\cTwoValue{0}$ \\
\cTwoOutcome{Paper} & $\cTwoMassHL{0.17}$ & $\cTwoValue{3}$ \\
\cTwoOutcome{Cardboard} & $\cTwoMassHL{0.15}$ & $\cTwoValue{3}$ \\
\cTwoOutcome{Plastic} & $\cTwoMassHL{0.25}$ & $\cTwoValue{8}$ \\
\cTwoOutcome{Steel} & $\cTwoMassHL{0.15}$ & $\cTwoValue{8}$ \\
\cTwoOutcome{Aluminium} & $\cTwoMassHL{0.2}$ & $\cTwoValue{14}$ \\
\bottomrule\end{tabular}
\end{center}
$\cTwoValue{S_X}=\cTwoValue{\{0,3,8,14\}}$; $p_X(x)=0$ elsewhere. The grouped PMF is:
\par\noindent\begin{minipage}[t]{.37\linewidth}\vspace{0pt}\centering
\begin{tabular}{cc}\toprule
$\cTwoValue{x}$ (\cTwoValue{USD}) & $\cTwoMassHL{p_X(x)}$ \\ \midrule
$\cTwoValue{0}$ & $\cTwoMassHL{0.08}$ \\
$\cTwoValue{3}$ & $\cTwoMassHL{0.32}$ \\
$\cTwoValue{8}$ & $\cTwoMassHL{0.4}$ \\
$\cTwoValue{14}$ & $\cTwoMassHL{0.2}$ \\
\bottomrule\end{tabular}
\end{minipage}\hfill\begin{minipage}[t]{.60\linewidth}\vspace{0pt}\centering
\GTenPMFGraph[color=cTwoMassColor!20,xhighlight=GTenHLValue,yhighlight=cTwoMassMarker,width=\linewidth,height=4.15cm,xlabel={$\cTwoValue{x}$ (\cTwoValue{USD})},ylabel={$\cTwoMassHL{p_X(x)}$},xstep=1]{0,3,8,14}{0.08,0.32,0.4,0.2}
\end{minipage}\par
\textbf{PMF check.} All masses are positive and $0.08+0.32+0.4+0.2=1$. The graph uses the actual numerical spacing; gaps carry no mass.
\task{2}{Accumulate in increasing support order: $0.08;\ 0.4;\ 0.8;\ 1$. Each new height is the previous height plus the next PMF mass.}
\par\noindent\begin{minipage}[t]{.37\linewidth}\vspace{0pt}\centering
\begin{tabular}{cc}\toprule
Cutoff interval & $\cTwoCDFHL{F_X(x)}$ \\ \midrule
$x<0$ & $\cTwoCDFHL{0}$ \\
$0\leq x<3$ & $\cTwoCDFHL{0.08}$ \\
$3\leq x<8$ & $\cTwoCDFHL{0.4}$ \\
$8\leq x<14$ & $\cTwoCDFHL{0.8}$ \\
$x\geq 14$ & $\cTwoCDFHL{1}$ \\
\bottomrule\end{tabular}
\end{minipage}\hfill\begin{minipage}[t]{.60\linewidth}\vspace{0pt}\centering
\begin{tikzpicture}\begin{axis}[classCDF,xmin=-3,xmax=17,xtick={0,3,8,14}]
\addplot[cTwoCDFColor,thick,<-] coordinates {(-3,0)(0,0)};
\addplot[cTwoCDFColor,thick] coordinates {(0,0.08)(3,0.08)};
\addplot[cTwoCDFColor,thick] coordinates {(3,0.4)(8,0.4)};
\addplot[cTwoCDFColor,thick] coordinates {(8,0.8)(14,0.8)};
\addplot[cTwoCDFColor,thick,->] coordinates {(14,1)(17,1)};
\addplot[only marks,mark=*,cTwoCDFColor,mark size=1.8pt,mark options={fill=white}] coordinates {(0,0)(3,0.08)(8,0.4)(14,0.8)};
\addplot[only marks,mark=*,cTwoCDFColor,mark size=1.8pt] coordinates {(0,0.08)(3,0.4)(8,0.8)(14,1)};
\end{axis}\end{tikzpicture}
\end{minipage}\par
\textbf{CDF check.} All heights lie in $[0,1]$ and never decrease. Filled left/open right endpoints make the graph right-continuous. Horizontal tails give limits 0 and 1; the CDF reaches 1 at 14 USD. The jumps, $0.08,\ 0.32,\ 0.4,\ 0.2$, reproduce the PMF.
\textbf{Interpretation.} $F_X(8)=0.8$: the probability that the buyback payment is at most 8 USD (every type except aluminium) is 0.8.
\vfill
\clearpage
\problemtitle{19}{Delivery loss claims}{D}
\textbf{Missing probability and capped payout.} A delivery insurer can receive only 0, 1, 2, 3 or 4 claims in one day. Their probabilities, in that order, are .15, .20, a missing probability, .20, .10. It pays 5 USD per claim, capped at 15 USD in total for the day; it pays the smaller of 5 times the count and 15 USD. There are no other payments. Record the daily payout.
\task{1}{Let $\cTwoValue{X}$ be the total payout for one day, in USD. The missing probability is $1-(.15+.20+.20+.10)=.35$. Apply $X=\min(5N,15)$ to the claim count $N$: payouts are $0,5,10,15,15$. Four claims would pay 20 USD without the cap. Group the last two categories.}
\begin{center}
\begin{tabular}{ccc}\toprule
Complete outcome & Probability & $\cTwoValue{X}$ (\cTwoValue{USD}) \\ \midrule
\cTwoOutcome{0 claims} & $\cTwoMassHL{0.15}$ & $\cTwoValue{0}$ \\
\cTwoOutcome{1 claims} & $\cTwoMassHL{0.2}$ & $\cTwoValue{5}$ \\
\cTwoOutcome{2 claims} & $\cTwoMassHL{0.35}$ & $\cTwoValue{10}$ \\
\cTwoOutcome{3 claims} & $\cTwoMassHL{0.2}$ & $\cTwoValue{15}$ \\
\cTwoOutcome{4 claims} & $\cTwoMassHL{0.1}$ & $\cTwoValue{15}$ \\
\bottomrule\end{tabular}
\end{center}
$\cTwoValue{S_X}=\cTwoValue{\{0,5,10,15\}}$; $p_X(x)=0$ elsewhere. The grouped PMF is:
\par\noindent\begin{minipage}[t]{.37\linewidth}\vspace{0pt}\centering
\begin{tabular}{cc}\toprule
$\cTwoValue{x}$ (\cTwoValue{USD}) & $\cTwoMassHL{p_X(x)}$ \\ \midrule
$\cTwoValue{0}$ & $\cTwoMassHL{0.15}$ \\
$\cTwoValue{5}$ & $\cTwoMassHL{0.2}$ \\
$\cTwoValue{10}$ & $\cTwoMassHL{0.35}$ \\
$\cTwoValue{15}$ & $\cTwoMassHL{0.3}$ \\
\bottomrule\end{tabular}
\end{minipage}\hfill\begin{minipage}[t]{.60\linewidth}\vspace{0pt}\centering
\GTenPMFGraph[color=cTwoMassColor!20,xhighlight=GTenHLValue,yhighlight=cTwoMassMarker,width=\linewidth,height=4.15cm,xlabel={$\cTwoValue{x}$ (\cTwoValue{USD})},ylabel={$\cTwoMassHL{p_X(x)}$},xstep=1]{0,5,10,15}{0.15,0.2,0.35,0.3}
\end{minipage}\par
\textbf{PMF check.} All masses are positive and $0.15+0.2+0.35+0.3=1$. The graph uses the actual numerical spacing; gaps carry no mass.
\task{2}{Accumulate in increasing support order: $0.15;\ 0.35;\ 0.7;\ 1$. Each new height is the previous height plus the next PMF mass.}
\par\noindent\begin{minipage}[t]{.37\linewidth}\vspace{0pt}\centering
\begin{tabular}{cc}\toprule
Cutoff interval & $\cTwoCDFHL{F_X(x)}$ \\ \midrule
$x<0$ & $\cTwoCDFHL{0}$ \\
$0\leq x<5$ & $\cTwoCDFHL{0.15}$ \\
$5\leq x<10$ & $\cTwoCDFHL{0.35}$ \\
$10\leq x<15$ & $\cTwoCDFHL{0.7}$ \\
$x\geq 15$ & $\cTwoCDFHL{1}$ \\
\bottomrule\end{tabular}
\end{minipage}\hfill\begin{minipage}[t]{.60\linewidth}\vspace{0pt}\centering
\begin{tikzpicture}\begin{axis}[classCDF,xmin=-5,xmax=20,xtick={0,5,10,15}]
\addplot[cTwoCDFColor,thick,<-] coordinates {(-5,0)(0,0)};
\addplot[cTwoCDFColor,thick] coordinates {(0,0.15)(5,0.15)};
\addplot[cTwoCDFColor,thick] coordinates {(5,0.35)(10,0.35)};
\addplot[cTwoCDFColor,thick] coordinates {(10,0.7)(15,0.7)};
\addplot[cTwoCDFColor,thick,->] coordinates {(15,1)(20,1)};
\addplot[only marks,mark=*,cTwoCDFColor,mark size=1.8pt,mark options={fill=white}] coordinates {(0,0)(5,0.15)(10,0.35)(15,0.7)};
\addplot[only marks,mark=*,cTwoCDFColor,mark size=1.8pt] coordinates {(0,0.15)(5,0.35)(10,0.7)(15,1)};
\end{axis}\end{tikzpicture}
\end{minipage}\par
\textbf{CDF check.} All heights lie in $[0,1]$ and never decrease. Filled left/open right endpoints make the graph right-continuous. Horizontal tails give limits 0 and 1; the CDF reaches 1 at 15 USD. The jumps, $0.15,\ 0.2,\ 0.35,\ 0.3$, reproduce the PMF.
\textbf{Interpretation.} $F_X(10)=0.7$: the probability that the payout is at most 10 USD (at most two claims) is 0.7.
\vfill
\clearpage
\problemtitle{20}{Event cancellation claims}{D}
\textbf{Missing probability and capped payout.} An event protection fund can receive only 0, 1, 2, 3 or 4 claims in one day. Their probabilities, in that order, are .22, .18, a missing probability, .17, .15. It pays 6 USD per claim, capped at 18 USD in total for the day; it pays the smaller of 6 times the count and 18 USD. There are no other payments. Record the daily payout.
\task{1}{Let $\cTwoValue{X}$ be the total payout for one day, in USD. The missing probability is $1-(.22+.18+.17+.15)=.28$. Apply $X=\min(6N,18)$ to the claim count $N$: payouts are $0,6,12,18,18$. Four claims would pay 24 USD without the cap. Group the last two categories.}
\begin{center}
\begin{tabular}{ccc}\toprule
Complete outcome & Probability & $\cTwoValue{X}$ (\cTwoValue{USD}) \\ \midrule
\cTwoOutcome{0 claims} & $\cTwoMassHL{0.22}$ & $\cTwoValue{0}$ \\
\cTwoOutcome{1 claims} & $\cTwoMassHL{0.18}$ & $\cTwoValue{6}$ \\
\cTwoOutcome{2 claims} & $\cTwoMassHL{0.28}$ & $\cTwoValue{12}$ \\
\cTwoOutcome{3 claims} & $\cTwoMassHL{0.17}$ & $\cTwoValue{18}$ \\
\cTwoOutcome{4 claims} & $\cTwoMassHL{0.15}$ & $\cTwoValue{18}$ \\
\bottomrule\end{tabular}
\end{center}
$\cTwoValue{S_X}=\cTwoValue{\{0,6,12,18\}}$; $p_X(x)=0$ elsewhere. The grouped PMF is:
\par\noindent\begin{minipage}[t]{.37\linewidth}\vspace{0pt}\centering
\begin{tabular}{cc}\toprule
$\cTwoValue{x}$ (\cTwoValue{USD}) & $\cTwoMassHL{p_X(x)}$ \\ \midrule
$\cTwoValue{0}$ & $\cTwoMassHL{0.22}$ \\
$\cTwoValue{6}$ & $\cTwoMassHL{0.18}$ \\
$\cTwoValue{12}$ & $\cTwoMassHL{0.28}$ \\
$\cTwoValue{18}$ & $\cTwoMassHL{0.32}$ \\
\bottomrule\end{tabular}
\end{minipage}\hfill\begin{minipage}[t]{.60\linewidth}\vspace{0pt}\centering
\GTenPMFGraph[color=cTwoMassColor!20,xhighlight=GTenHLValue,yhighlight=cTwoMassMarker,width=\linewidth,height=4.15cm,xlabel={$\cTwoValue{x}$ (\cTwoValue{USD})},ylabel={$\cTwoMassHL{p_X(x)}$},xstep=1]{0,6,12,18}{0.22,0.18,0.28,0.32}
\end{minipage}\par
\textbf{PMF check.} All masses are positive and $0.22+0.18+0.28+0.32=1$. The graph uses the actual numerical spacing; gaps carry no mass.
\task{2}{Accumulate in increasing support order: $0.22;\ 0.4;\ 0.68;\ 1$. Each new height is the previous height plus the next PMF mass.}
\par\noindent\begin{minipage}[t]{.37\linewidth}\vspace{0pt}\centering
\begin{tabular}{cc}\toprule
Cutoff interval & $\cTwoCDFHL{F_X(x)}$ \\ \midrule
$x<0$ & $\cTwoCDFHL{0}$ \\
$0\leq x<6$ & $\cTwoCDFHL{0.22}$ \\
$6\leq x<12$ & $\cTwoCDFHL{0.4}$ \\
$12\leq x<18$ & $\cTwoCDFHL{0.68}$ \\
$x\geq 18$ & $\cTwoCDFHL{1}$ \\
\bottomrule\end{tabular}
\end{minipage}\hfill\begin{minipage}[t]{.60\linewidth}\vspace{0pt}\centering
\begin{tikzpicture}\begin{axis}[classCDF,xmin=-6,xmax=24,xtick={0,6,12,18}]
\addplot[cTwoCDFColor,thick,<-] coordinates {(-6,0)(0,0)};
\addplot[cTwoCDFColor,thick] coordinates {(0,0.22)(6,0.22)};
\addplot[cTwoCDFColor,thick] coordinates {(6,0.4)(12,0.4)};
\addplot[cTwoCDFColor,thick] coordinates {(12,0.68)(18,0.68)};
\addplot[cTwoCDFColor,thick,->] coordinates {(18,1)(24,1)};
\addplot[only marks,mark=*,cTwoCDFColor,mark size=1.8pt,mark options={fill=white}] coordinates {(0,0)(6,0.22)(12,0.4)(18,0.68)};
\addplot[only marks,mark=*,cTwoCDFColor,mark size=1.8pt] coordinates {(0,0.22)(6,0.4)(12,0.68)(18,1)};
\end{axis}\end{tikzpicture}
\end{minipage}\par
\textbf{CDF check.} All heights lie in $[0,1]$ and never decrease. Filled left/open right endpoints make the graph right-continuous. Horizontal tails give limits 0 and 1; the CDF reaches 1 at 18 USD. The jumps, $0.22,\ 0.18,\ 0.28,\ 0.32$, reproduce the PMF.
\textbf{Interpretation.} $F_X(12)=0.68$: the probability that the payout is at most 12 USD (at most two claims) is 0.68.
\vfill
\clearpage
\problemtitle{21}{Tote bag sales}{E}
\textbf{Independent purchases and discount.} A school enterprise offers one tote bag to each of three visitors. Independently, each buys one with probability 0.3 or passes with probability 0.7, regardless of the discount. Each item costs 7 USD. Only if all three buy, subtract a 5 USD discount from the combined bill. Otherwise there is no discount. There are no other sales or charges. Record total revenue.
\task{1}{Let $\cTwoValue{X}$ be the total revenue from the three visitors, in USD. Use B for buying and N for passing, in visitor order. Multiply within each complete sequence: e.g. $P(\mathrm{BNN})=0.3(0.7)^2=0.147$. The three one-buyer and three two-buyer outcomes each form a group. Only BBB receives the discount: $3(7)-5=16$ USD.}
\begin{center}
\begin{tabular}{ccc}\toprule
Complete outcome & Probability & $\cTwoValue{X}$ (\cTwoValue{USD}) \\ \midrule
\cTwoOutcome{NNN} & $\cTwoMassHL{0.343}$ & $\cTwoValue{0}$ \\
\cTwoOutcome{BNN} & $\cTwoMassHL{0.147}$ & $\cTwoValue{7}$ \\
\cTwoOutcome{NBN} & $\cTwoMassHL{0.147}$ & $\cTwoValue{7}$ \\
\cTwoOutcome{NNB} & $\cTwoMassHL{0.147}$ & $\cTwoValue{7}$ \\
\cTwoOutcome{BBN} & $\cTwoMassHL{0.063}$ & $\cTwoValue{14}$ \\
\cTwoOutcome{BNB} & $\cTwoMassHL{0.063}$ & $\cTwoValue{14}$ \\
\cTwoOutcome{NBB} & $\cTwoMassHL{0.063}$ & $\cTwoValue{14}$ \\
\cTwoOutcome{BBB} & $\cTwoMassHL{0.027}$ & $\cTwoValue{16}$ \\
\bottomrule\end{tabular}
\end{center}
$\cTwoValue{S_X}=\cTwoValue{\{0,7,14,16\}}$; $p_X(x)=0$ elsewhere. The grouped PMF is:
\par\noindent\begin{minipage}[t]{.37\linewidth}\vspace{0pt}\centering
\begin{tabular}{cc}\toprule
$\cTwoValue{x}$ (\cTwoValue{USD}) & $\cTwoMassHL{p_X(x)}$ \\ \midrule
$\cTwoValue{0}$ & $\cTwoMassHL{0.343}$ \\
$\cTwoValue{7}$ & $\cTwoMassHL{0.441}$ \\
$\cTwoValue{14}$ & $\cTwoMassHL{0.189}$ \\
$\cTwoValue{16}$ & $\cTwoMassHL{0.027}$ \\
\bottomrule\end{tabular}
\end{minipage}\hfill\begin{minipage}[t]{.60\linewidth}\vspace{0pt}\centering
\GTenPMFGraph[color=cTwoMassColor!20,xhighlight=GTenHLValue,yhighlight=cTwoMassMarker,width=\linewidth,height=4.15cm,xlabel={$\cTwoValue{x}$ (\cTwoValue{USD})},ylabel={$\cTwoMassHL{p_X(x)}$},xstep=1]{0,7,14,16}{0.343,0.441,0.189,0.027}
\end{minipage}\par
\textbf{PMF check.} All masses are positive and $0.343+0.441+0.189+0.027=1$. The graph uses the actual numerical spacing; gaps carry no mass.
\task{2}{Accumulate in increasing support order: $0.343;\ 0.784;\ 0.973;\ 1$. Each new height is the previous height plus the next PMF mass.}
\par\noindent\begin{minipage}[t]{.37\linewidth}\vspace{0pt}\centering
\begin{tabular}{cc}\toprule
Cutoff interval & $\cTwoCDFHL{F_X(x)}$ \\ \midrule
$x<0$ & $\cTwoCDFHL{0}$ \\
$0\leq x<7$ & $\cTwoCDFHL{0.343}$ \\
$7\leq x<14$ & $\cTwoCDFHL{0.784}$ \\
$14\leq x<16$ & $\cTwoCDFHL{0.973}$ \\
$x\geq 16$ & $\cTwoCDFHL{1}$ \\
\bottomrule\end{tabular}
\end{minipage}\hfill\begin{minipage}[t]{.60\linewidth}\vspace{0pt}\centering
\begin{tikzpicture}\begin{axis}[classCDF,xmin=-2,xmax=18,xtick={0,7,14,16}]
\addplot[cTwoCDFColor,thick,<-] coordinates {(-2,0)(0,0)};
\addplot[cTwoCDFColor,thick] coordinates {(0,0.343)(7,0.343)};
\addplot[cTwoCDFColor,thick] coordinates {(7,0.784)(14,0.784)};
\addplot[cTwoCDFColor,thick] coordinates {(14,0.973)(16,0.973)};
\addplot[cTwoCDFColor,thick,->] coordinates {(16,1)(18,1)};
\addplot[only marks,mark=*,cTwoCDFColor,mark size=1.8pt,mark options={fill=white}] coordinates {(0,0)(7,0.343)(14,0.784)(16,0.973)};
\addplot[only marks,mark=*,cTwoCDFColor,mark size=1.8pt] coordinates {(0,0.343)(7,0.784)(14,0.973)(16,1)};
\end{axis}\end{tikzpicture}
\end{minipage}\par
\textbf{CDF check.} All heights lie in $[0,1]$ and never decrease. Filled left/open right endpoints make the graph right-continuous. Horizontal tails give limits 0 and 1; the CDF reaches 1 at 16 USD. The jumps, $0.343,\ 0.441,\ 0.189,\ 0.027$, reproduce the PMF.
\textbf{Interpretation.} $F_X(14)=0.973$: the probability that revenue is at most 14 USD (not all three visitors buy) is 0.973.
\vfill
\clearpage
\problemtitle{22}{Poster sales}{E}
\textbf{Independent purchases and discount.} A school enterprise offers one poster to each of three visitors. Independently, each buys one with probability 0.7 or passes with probability 0.3, regardless of the discount. Each item costs 8 USD. Only if all three buy, subtract a 3 USD discount from the combined bill. Otherwise there is no discount. There are no other sales or charges. Record total revenue.
\task{1}{Let $\cTwoValue{X}$ be the total revenue from the three visitors, in USD. Use B for buying and N for passing, in visitor order. Multiply within each complete sequence: e.g. $P(\mathrm{BNN})=0.7(0.3)^2=0.063$. The three one-buyer and three two-buyer outcomes each form a group. Only BBB receives the discount: $3(8)-3=21$ USD.}
\begin{center}
\begin{tabular}{ccc}\toprule
Complete outcome & Probability & $\cTwoValue{X}$ (\cTwoValue{USD}) \\ \midrule
\cTwoOutcome{NNN} & $\cTwoMassHL{0.027}$ & $\cTwoValue{0}$ \\
\cTwoOutcome{BNN} & $\cTwoMassHL{0.063}$ & $\cTwoValue{8}$ \\
\cTwoOutcome{NBN} & $\cTwoMassHL{0.063}$ & $\cTwoValue{8}$ \\
\cTwoOutcome{NNB} & $\cTwoMassHL{0.063}$ & $\cTwoValue{8}$ \\
\cTwoOutcome{BBN} & $\cTwoMassHL{0.147}$ & $\cTwoValue{16}$ \\
\cTwoOutcome{BNB} & $\cTwoMassHL{0.147}$ & $\cTwoValue{16}$ \\
\cTwoOutcome{NBB} & $\cTwoMassHL{0.147}$ & $\cTwoValue{16}$ \\
\cTwoOutcome{BBB} & $\cTwoMassHL{0.343}$ & $\cTwoValue{21}$ \\
\bottomrule\end{tabular}
\end{center}
$\cTwoValue{S_X}=\cTwoValue{\{0,8,16,21\}}$; $p_X(x)=0$ elsewhere. The grouped PMF is:
\par\noindent\begin{minipage}[t]{.37\linewidth}\vspace{0pt}\centering
\begin{tabular}{cc}\toprule
$\cTwoValue{x}$ (\cTwoValue{USD}) & $\cTwoMassHL{p_X(x)}$ \\ \midrule
$\cTwoValue{0}$ & $\cTwoMassHL{0.027}$ \\
$\cTwoValue{8}$ & $\cTwoMassHL{0.189}$ \\
$\cTwoValue{16}$ & $\cTwoMassHL{0.441}$ \\
$\cTwoValue{21}$ & $\cTwoMassHL{0.343}$ \\
\bottomrule\end{tabular}
\end{minipage}\hfill\begin{minipage}[t]{.60\linewidth}\vspace{0pt}\centering
\GTenPMFGraph[color=cTwoMassColor!20,xhighlight=GTenHLValue,yhighlight=cTwoMassMarker,width=\linewidth,height=4.15cm,xlabel={$\cTwoValue{x}$ (\cTwoValue{USD})},ylabel={$\cTwoMassHL{p_X(x)}$},xstep=1]{0,8,16,21}{0.027,0.189,0.441,0.343}
\end{minipage}\par
\textbf{PMF check.} All masses are positive and $0.027+0.189+0.441+0.343=1$. The graph uses the actual numerical spacing; gaps carry no mass.
\task{2}{Accumulate in increasing support order: $0.027;\ 0.216;\ 0.657;\ 1$. Each new height is the previous height plus the next PMF mass.}
\par\noindent\begin{minipage}[t]{.37\linewidth}\vspace{0pt}\centering
\begin{tabular}{cc}\toprule
Cutoff interval & $\cTwoCDFHL{F_X(x)}$ \\ \midrule
$x<0$ & $\cTwoCDFHL{0}$ \\
$0\leq x<8$ & $\cTwoCDFHL{0.027}$ \\
$8\leq x<16$ & $\cTwoCDFHL{0.216}$ \\
$16\leq x<21$ & $\cTwoCDFHL{0.657}$ \\
$x\geq 21$ & $\cTwoCDFHL{1}$ \\
\bottomrule\end{tabular}
\end{minipage}\hfill\begin{minipage}[t]{.60\linewidth}\vspace{0pt}\centering
\begin{tikzpicture}\begin{axis}[classCDF,xmin=-5,xmax=26,xtick={0,8,16,21}]
\addplot[cTwoCDFColor,thick,<-] coordinates {(-5,0)(0,0)};
\addplot[cTwoCDFColor,thick] coordinates {(0,0.027)(8,0.027)};
\addplot[cTwoCDFColor,thick] coordinates {(8,0.216)(16,0.216)};
\addplot[cTwoCDFColor,thick] coordinates {(16,0.657)(21,0.657)};
\addplot[cTwoCDFColor,thick,->] coordinates {(21,1)(26,1)};
\addplot[only marks,mark=*,cTwoCDFColor,mark size=1.8pt,mark options={fill=white}] coordinates {(0,0)(8,0.027)(16,0.216)(21,0.657)};
\addplot[only marks,mark=*,cTwoCDFColor,mark size=1.8pt] coordinates {(0,0.027)(8,0.216)(16,0.657)(21,1)};
\end{axis}\end{tikzpicture}
\end{minipage}\par
\textbf{CDF check.} All heights lie in $[0,1]$ and never decrease. Filled left/open right endpoints make the graph right-continuous. Horizontal tails give limits 0 and 1; the CDF reaches 1 at 21 USD. The jumps, $0.027,\ 0.189,\ 0.441,\ 0.343$, reproduce the PMF.
\textbf{Interpretation.} $F_X(16)=0.657$: the probability that revenue is at most 16 USD (not all three visitors buy) is 0.657.
\vfill
\clearpage
\problemtitle{23}{Two-stage media fund}{F}
\textbf{Signed two-stage net results.} A hypothetical classroom fund receives two cash returns. The first is 17 USD, 0 USD or 7 USD with probabilities .20, .30, .50, respectively. Independently, the second is 0 USD with probability 0.4 or 7 USD with probability 0.6. These are all possible returns. A fixed participation cost of 7 USD is charged once. There are no other receipts or charges. Record the sum of both returns minus the cost; a negative net result means a loss.
\task{1}{Let $\cTwoValue{X}$ be the signed net result after both stages and the cost, in USD. For each ordered pair, multiply independent probabilities, then compute $X=R_1+R_2-7$. For example, $(17,0)$ has probability $.20(0.4)=0.08$ and net result $17+0-7=10$. The pairs $(0,7)$ and $(7,0)$ both break even; group their masses. Subtract the cost only once and sort signed results numerically.}
\begin{center}
\begin{tabular}{ccc}\toprule
Complete outcome & Probability & $\cTwoValue{X}$ (\cTwoValue{USD}) \\ \midrule
\cTwoOutcome{$(17,0)$} & $\cTwoMassHL{0.08}$ & $\cTwoValue{10}$ \\
\cTwoOutcome{$(17,7)$} & $\cTwoMassHL{0.12}$ & $\cTwoValue{17}$ \\
\cTwoOutcome{$(0,0)$} & $\cTwoMassHL{0.12}$ & $\cTwoValue{-7}$ \\
\cTwoOutcome{$(0,7)$} & $\cTwoMassHL{0.18}$ & $\cTwoValue{0}$ \\
\cTwoOutcome{$(7,0)$} & $\cTwoMassHL{0.2}$ & $\cTwoValue{0}$ \\
\cTwoOutcome{$(7,7)$} & $\cTwoMassHL{0.3}$ & $\cTwoValue{7}$ \\
\bottomrule\end{tabular}
\end{center}
$\cTwoValue{S_X}=\cTwoValue{\{-7,0,7,10,17\}}$; $p_X(x)=0$ elsewhere. The grouped PMF is:
\par\noindent\begin{minipage}[t]{.37\linewidth}\vspace{0pt}\centering
\begin{tabular}{cc}\toprule
$\cTwoValue{x}$ (\cTwoValue{USD}) & $\cTwoMassHL{p_X(x)}$ \\ \midrule
$\cTwoValue{-7}$ & $\cTwoMassHL{0.12}$ \\
$\cTwoValue{0}$ & $\cTwoMassHL{0.38}$ \\
$\cTwoValue{7}$ & $\cTwoMassHL{0.3}$ \\
$\cTwoValue{10}$ & $\cTwoMassHL{0.08}$ \\
$\cTwoValue{17}$ & $\cTwoMassHL{0.12}$ \\
\bottomrule\end{tabular}
\end{minipage}\hfill\begin{minipage}[t]{.60\linewidth}\vspace{0pt}\centering
\GTenPMFGraph[color=cTwoMassColor!20,xhighlight=GTenHLValue,yhighlight=cTwoMassMarker,width=\linewidth,height=4.15cm,xlabel={$\cTwoValue{x}$ (\cTwoValue{USD})},ylabel={$\cTwoMassHL{p_X(x)}$},xstep=1]{-7,0,7,10,17}{0.12,0.38,0.3,0.08,0.12}
\end{minipage}\par
\textbf{PMF check.} All masses are positive and $0.12+0.38+0.3+0.08+0.12=1$. The graph uses the actual numerical spacing; gaps carry no mass.
\task{2}{Accumulate in increasing support order: $0.12;\ 0.5;\ 0.8;\ 0.88;\ 1$. Each new height is the previous height plus the next PMF mass.}
\par\noindent\begin{minipage}[t]{.37\linewidth}\vspace{0pt}\centering
\begin{tabular}{cc}\toprule
Cutoff interval & $\cTwoCDFHL{F_X(x)}$ \\ \midrule
$x<-7$ & $\cTwoCDFHL{0}$ \\
$-7\leq x<0$ & $\cTwoCDFHL{0.12}$ \\
$0\leq x<7$ & $\cTwoCDFHL{0.5}$ \\
$7\leq x<10$ & $\cTwoCDFHL{0.8}$ \\
$10\leq x<17$ & $\cTwoCDFHL{0.88}$ \\
$x\geq 17$ & $\cTwoCDFHL{1}$ \\
\bottomrule\end{tabular}
\end{minipage}\hfill\begin{minipage}[t]{.60\linewidth}\vspace{0pt}\centering
\begin{tikzpicture}\begin{axis}[classCDF,xmin=-10,xmax=20,xtick={-7,0,7,10,17}]
\addplot[cTwoCDFColor,thick,<-] coordinates {(-10,0)(-7,0)};
\addplot[cTwoCDFColor,thick] coordinates {(-7,0.12)(0,0.12)};
\addplot[cTwoCDFColor,thick] coordinates {(0,0.5)(7,0.5)};
\addplot[cTwoCDFColor,thick] coordinates {(7,0.8)(10,0.8)};
\addplot[cTwoCDFColor,thick] coordinates {(10,0.88)(17,0.88)};
\addplot[cTwoCDFColor,thick,->] coordinates {(17,1)(20,1)};
\addplot[only marks,mark=*,cTwoCDFColor,mark size=1.8pt,mark options={fill=white}] coordinates {(-7,0)(0,0.12)(7,0.5)(10,0.8)(17,0.88)};
\addplot[only marks,mark=*,cTwoCDFColor,mark size=1.8pt] coordinates {(-7,0.12)(0,0.5)(7,0.8)(10,0.88)(17,1)};
\end{axis}\end{tikzpicture}
\end{minipage}\par
\textbf{CDF check.} All heights lie in $[0,1]$ and never decrease. Filled left/open right endpoints make the graph right-continuous. Horizontal tails give limits 0 and 1; the CDF reaches 1 at 17 USD. The jumps, $0.12,\ 0.38,\ 0.3,\ 0.08,\ 0.12$, reproduce the PMF.
\textbf{Interpretation.} $F_X(0)=0.5$: the probability that the net result is nonpositive (a loss or break-even) is 0.5.
\vfill
\clearpage
\problemtitle{24}{Two-stage sports fund}{F}
\textbf{Signed two-stage net results.} A hypothetical classroom fund receives two cash returns. The first is 19 USD, 0 USD or 8 USD with probabilities .35, .15, .50, respectively. Independently, the second is 0 USD with probability 0.8 or 8 USD with probability 0.2. These are all possible returns. A fixed participation cost of 8 USD is charged once. There are no other receipts or charges. Record the sum of both returns minus the cost; a negative net result means a loss.
\task{1}{Let $\cTwoValue{X}$ be the signed net result after both stages and the cost, in USD. For each ordered pair, multiply independent probabilities, then compute $X=R_1+R_2-8$. For example, $(19,0)$ has probability $.35(0.8)=0.28$ and net result $19+0-8=11$. The pairs $(0,8)$ and $(8,0)$ both break even; group their masses. Subtract the cost only once and sort signed results numerically.}
\begin{center}
\begin{tabular}{ccc}\toprule
Complete outcome & Probability & $\cTwoValue{X}$ (\cTwoValue{USD}) \\ \midrule
\cTwoOutcome{$(19,0)$} & $\cTwoMassHL{0.28}$ & $\cTwoValue{11}$ \\
\cTwoOutcome{$(19,8)$} & $\cTwoMassHL{0.07}$ & $\cTwoValue{19}$ \\
\cTwoOutcome{$(0,0)$} & $\cTwoMassHL{0.12}$ & $\cTwoValue{-8}$ \\
\cTwoOutcome{$(0,8)$} & $\cTwoMassHL{0.03}$ & $\cTwoValue{0}$ \\
\cTwoOutcome{$(8,0)$} & $\cTwoMassHL{0.4}$ & $\cTwoValue{0}$ \\
\cTwoOutcome{$(8,8)$} & $\cTwoMassHL{0.1}$ & $\cTwoValue{8}$ \\
\bottomrule\end{tabular}
\end{center}
$\cTwoValue{S_X}=\cTwoValue{\{-8,0,8,11,19\}}$; $p_X(x)=0$ elsewhere. The grouped PMF is:
\par\noindent\begin{minipage}[t]{.37\linewidth}\vspace{0pt}\centering
\begin{tabular}{cc}\toprule
$\cTwoValue{x}$ (\cTwoValue{USD}) & $\cTwoMassHL{p_X(x)}$ \\ \midrule
$\cTwoValue{-8}$ & $\cTwoMassHL{0.12}$ \\
$\cTwoValue{0}$ & $\cTwoMassHL{0.43}$ \\
$\cTwoValue{8}$ & $\cTwoMassHL{0.1}$ \\
$\cTwoValue{11}$ & $\cTwoMassHL{0.28}$ \\
$\cTwoValue{19}$ & $\cTwoMassHL{0.07}$ \\
\bottomrule\end{tabular}
\end{minipage}\hfill\begin{minipage}[t]{.60\linewidth}\vspace{0pt}\centering
\GTenPMFGraph[color=cTwoMassColor!20,xhighlight=GTenHLValue,yhighlight=cTwoMassMarker,width=\linewidth,height=4.15cm,xlabel={$\cTwoValue{x}$ (\cTwoValue{USD})},ylabel={$\cTwoMassHL{p_X(x)}$},xstep=1]{-8,0,8,11,19}{0.12,0.43,0.1,0.28,0.07}
\end{minipage}\par
\textbf{PMF check.} All masses are positive and $0.12+0.43+0.1+0.28+0.07=1$. The graph uses the actual numerical spacing; gaps carry no mass.
\task{2}{Accumulate in increasing support order: $0.12;\ 0.55;\ 0.65;\ 0.93;\ 1$. Each new height is the previous height plus the next PMF mass.}
\par\noindent\begin{minipage}[t]{.37\linewidth}\vspace{0pt}\centering
\begin{tabular}{cc}\toprule
Cutoff interval & $\cTwoCDFHL{F_X(x)}$ \\ \midrule
$x<-8$ & $\cTwoCDFHL{0}$ \\
$-8\leq x<0$ & $\cTwoCDFHL{0.12}$ \\
$0\leq x<8$ & $\cTwoCDFHL{0.55}$ \\
$8\leq x<11$ & $\cTwoCDFHL{0.65}$ \\
$11\leq x<19$ & $\cTwoCDFHL{0.93}$ \\
$x\geq 19$ & $\cTwoCDFHL{1}$ \\
\bottomrule\end{tabular}
\end{minipage}\hfill\begin{minipage}[t]{.60\linewidth}\vspace{0pt}\centering
\begin{tikzpicture}\begin{axis}[classCDF,xmin=-11,xmax=22,xtick={-8,0,8,11,19}]
\addplot[cTwoCDFColor,thick,<-] coordinates {(-11,0)(-8,0)};
\addplot[cTwoCDFColor,thick] coordinates {(-8,0.12)(0,0.12)};
\addplot[cTwoCDFColor,thick] coordinates {(0,0.55)(8,0.55)};
\addplot[cTwoCDFColor,thick] coordinates {(8,0.65)(11,0.65)};
\addplot[cTwoCDFColor,thick] coordinates {(11,0.93)(19,0.93)};
\addplot[cTwoCDFColor,thick,->] coordinates {(19,1)(22,1)};
\addplot[only marks,mark=*,cTwoCDFColor,mark size=1.8pt,mark options={fill=white}] coordinates {(-8,0)(0,0.12)(8,0.55)(11,0.65)(19,0.93)};
\addplot[only marks,mark=*,cTwoCDFColor,mark size=1.8pt] coordinates {(-8,0.12)(0,0.55)(8,0.65)(11,0.93)(19,1)};
\end{axis}\end{tikzpicture}
\end{minipage}\par
\textbf{CDF check.} All heights lie in $[0,1]$ and never decrease. Filled left/open right endpoints make the graph right-continuous. Horizontal tails give limits 0 and 1; the CDF reaches 1 at 19 USD. The jumps, $0.12,\ 0.43,\ 0.1,\ 0.28,\ 0.07$, reproduce the PMF.
\textbf{Interpretation.} $F_X(0)=0.55$: the probability that the net result is nonpositive (a loss or break-even) is 0.55.
\vfill
\fi
\end{document}
